% Generated by assemble.py; edit main.tex and its inputs.
\documentclass[11pt]{article}
\usepackage[T1]{fontenc}
\usepackage{lmodern,amsmath,amssymb,amsthm,mathtools}
\usepackage[margin=1in]{geometry}
\usepackage{booktabs,longtable,array,microtype,xcolor}
\usepackage{xurl,hyperref,fancyhdr}
\hypersetup{colorlinks=true,linkcolor=blue!45!black,urlcolor=blue!45!black,citecolor=blue!45!black}
\hypersetup{pdftitle={Polynomial Centers for Diophantine Consistency},
pdfauthor={AI-assisted research prepared for Vladimir Reshetnikov and ProveIt},
pdfsubject={Degree 208 in ten positive variables and direct quartic companions},
pdfkeywords={Diophantine equation, Con(ZFC), polynomial center, Pell equation}}
\setlength{\emergencystretch}{3em}
\allowdisplaybreaks[4]
\newtheorem{theorem}{Theorem}[section]
\newtheorem{lemma}[theorem]{Lemma}
\newtheorem{proposition}[theorem]{Proposition}
\newtheorem{corollary}[theorem]{Corollary}
\theoremstyle{remark}\newtheorem{remark}[theorem]{Remark}
\newcommand{\Z}{\mathbb Z}\newcommand{\N}{\mathbb N}
\newcommand{\Pos}{\mathbb Z_{>0}}
\newcommand{\Con}{\operatorname{Con}}
\newcommand{\ZFC}{\mathsf{ZFC}}\newcommand{\ZF}{\mathsf{ZF}}
\newcommand{\topdeg}[2]{[#1]_{#2}}
\newcommand{\qpoly}{\operatorname{quo}}
\newcommand{\code}[1]{\texttt{\detokenize{#1}}}
% BEGIN embedded metrics.tex
% Generated by generate_inputs.py from whole_product/equation.json.
\newcommand{\PrimaryOps}{116}
\newcommand{\PrimaryMul}{71}
\newcommand{\PrimaryAdd}{24}
\newcommand{\PrimarySub}{21}

% END embedded metrics.tex

\pagestyle{fancy}\fancyhf{}
\fancyhead[L]{\small Polynomial centers for Diophantine consistency}
\fancyhead[R]{\small 25 October 2026}
\fancyfoot[C]{\thepage}
\setlength{\headheight}{14pt}
\title{Polynomial Centers for Diophantine Consistency\\[0.4em]
\large A degree-208 equation in ten positive variables,\\
with direct quartic companions in 203 naturals and 233 integers}
\author{AI-assisted research prepared for Vladimir Reshetnikov and ProveIt}
\date{25 October 2026}
\begin{document}
\maketitle
\begin{abstract}
We give a literal integer polynomial in ten strictly positive variables,
of exact total degree 208, evaluated by \PrimaryOps{} live binary arithmetic
operations. Its unsolvability is equivalent to $\Con(\ZFC)$ by an ordinary
proof formalizable in ZFC, composed with explicitly identified inherited
arithmetic and refutation calibrations. The new step centers the whole Pell
product $gV$, using polynomial division of $EV$ and a truncated square-root
expansion. All divisions disappear from the final formula. A signed
majorant argument establishes soundness before any positivity conclusion
about the encoded computation is used. This also resolves the previously
unproved square-transport case in source 163.

The package includes smaller-operation alternatives at degree 232, 228,
224 and 220, with the same ten positive variables and the same entire
fixed coefficient tuple. A separate direct-calculus component gives fully
expanded quartics in 203 nonnegative variables and 233 unrestricted
integers. The direct component does not simulate K215 or use Jones's
universal equation. No global minimum, proof of consistency, or new
end-to-end proof-assistant certificate is claimed.
\end{abstract}
\clearpage
\setcounter{tocdepth}{1}
\tableofcontents
\clearpage

\section{Result and comparison boundary}

The objective is an actual polynomial with fixed integer coefficients,
not an exponential Diophantine formula or an unspecified instance of MRDP.
Throughout, $\N=\{0,1,2,\ldots\}$ and $\Pos=\{1,2,3,\ldots\}$.
The domain is part of every assertion.

\begin{theorem}[Main equation]\label{thm:main}
Let $P_{208}$ be the polynomial defined in Sections~\ref{sec:formula}
and~\ref{sec:center}, with the four fixed integers defined in
Appendix~\ref{app:constants}. Then $P_{208}\in\Z[J,F,\alpha,i,T,w,s,n,\eta,\zeta]$,
its exact total degree is 208, and
\[
 \ZFC\ \vdash\
 \bigl(\neg\exists\mathbf x\in\Pos^{10}\;P_{208}(\mathbf x)=0\bigr)
 \ \longleftrightarrow\ \Con(\ZFC).
\]
The displayed provability claim is established by the ordinary proofs here and
the precisely named inherited interfaces. This package does not contain
a newly checked object-language ZFC derivation or a newly compiled
end-to-end Lean or Rocq development.
\end{theorem}

We use the immutable repository snapshot
\begin{center}\small\ttfamily
fb1bed5eb7196cf35abfaa746c18c6ee7eae6cfb.
\end{center}
The comparison includes recent packages 163--165, not merely the older
summary table. Source 163 gives ten-positive degree 232 in 87 operations;
source 164 gives degree 304 in 85 operations~\cite{p163,p164}. Their two mechanisms compose
to give 85/232. The further centers below give new degree/cost choices.
\begin{center}
\begin{tabular}{lrrl}
\toprule
Equation & Operations & Exact degree & Main change\\
\midrule
Source 164 &85&304&Inverse Lucas correction\\
Source 163 &87&232&Unsquared strong trace\\
Trace plus inverse &85&232&Composition of the two\\
Lower center &87&228&$(A-1)UE$\\
Center at $AUE$ &88&224&First norm cancellation\\
Quotient center &96&220&Divide $E$ by $A$\\
\textbf{Whole-product center} &\textbf{\PrimaryOps}&\textbf{208}&\textbf{Divide $EV$ by powers of $A$}\\
\bottomrule
\end{tabular}
\end{center}
All these counts concern ten independent positive variables and the same
whole coefficient tuple. Every live binary addition, subtraction and
multiplication is charged, including squares and fixed scalar products.
Named values are polynomial abbreviations, not extra unknowns. These
operation counts are not bit-complexity estimates: the largest fixed
coefficient port has 592,884 bits.

\begin{theorem}[Direct quartic companions]\label{thm:direct}
The literal residual lists in the package define quartics
\[
 P_{203}^{\N}=\sum_{j=0}^{154}q_j^2,
 \qquad P_{233}^{\Z}=\sum_{j=0}^{168}r_j^2,
\]
in respectively 203 nonnegative and 233 unrestricted integer variables,
such that ZFC proves each no-zero statement equivalent to $\Con(\ZFC)$.
All $q_j,r_j$ are explicitly listed integer quadratics; complete collected
quartic coefficients are also supplied. These constructions use the
numerical proof calculus C222 and its ZF/ZFC calibration.
\end{theorem}

This is a set of measured improvements, not a claim that one ordering of
``simplicity'' dominates every other. The repository already has fewer-variable
equations of much larger degree. Coefficient size, monomial count, proof
burden, variable domain, degree and evaluation cost must be kept separate.

The ten-positive-variable equation immediately becomes a ten-nonnegative-variable
equation of the same degree by replacing each input $x_j$ with $u_j+1$.
If an unrestricted integer domain is required, the four-square theorem gives
the completely explicit alternative
\[
 P_{208}\left(1+\sum_{r=1}^4z_{1r}^2,\ldots,
                   1+\sum_{r=1}^4z_{10,r}^2\right)=0,
 \qquad \mathbf z\in\Z^{40},
\]
of exact degree 416. Each positive input has such a representation.
The degree is exact because restricting each block to one square leaves
the nonzero leading polynomial with every exponent doubled. This standard
domain conversion is separate from the direct quartic reductions.

\section{The complete common polynomial definitions}\label{sec:formula}

Quantify exactly
\[
 J,F,\alpha,i,T,w,s,n,\eta,\zeta\in\Pos.
\]
The fixed integers are $\beta,K_0,M_C,M_F$. They are supplied both as
authenticated binary numerals and by the complete finite recipe in
Appendix~\ref{app:constants}. In particular,
\[
 \beta=2^{590069}-1,\qquad
 0<M_C<\beta<M_F<2\beta,\qquad K_0>0.
\]
These inequalities alone do not calibrate an arbitrary four-tuple. The
construction uses all four actual integers together.

Define the following ordinary polynomials in the ten variables:
\begin{align*}
 h&=\beta J,&q&=h+1,&C&=q-F-\alpha,\\
 X&=qw,&Y&=qs,&k&=\eta+\zeta,\\
 c&=kY+\eta,&A&=Y(X+1)+1,&\Delta&=A^2-1,\\
 H&=4A-5,&\ell&=(XY)(kY),&E&=1+\ell(\ell+k),\\
 R&=[h(q-F)+\alpha](q^2-1)+(M_C+qM_F)J,\\
 B&=(K_0+w)C-F,&U&=ic^2,&a&=\Delta U^2E,\\
 F_0&=1+a,&G&=EF_0,&f&=1+2a,\\
 Q&=4\Delta aF_0,&b&=cT-R,&V&=fb-c,\\
 K&=X(3c-2X)+2,&D&=6q^6,&M&=hDH,\\
 e_6&=6(k-R-1)-\ell(k^2-4).
\end{align*}
The letters $C,F,F_0,G,H,K$ here have only the meanings displayed above.
They are not coefficients to be chosen or additional quantifiers.

For any polynomial shift $\Omega$ and positive fixed integer $r$, put
\begin{align}
 Z_0&=hD\{H(n-b^2B)+K\}-H\{DB+h(e_6-6q^4s)\},\label{eq:Z0}\\
 Z&=Z_0-M\Omega,\label{eq:Z}\\
 P_{\Omega,r}
 &=Q\bigl(Z^2-r^2M^2GV^2\bigr)+\bigl(r^2M^2G-Z^2\bigr).
 \label{eq:generalP}
\end{align}
The principal equation uses $r=4$ and the explicitly polynomial shift in
the next section. Equivalently, its literal evaluator inserts
$n-b^2B-\Omega$ inside the braces in~\eqref{eq:Z0}, saving one operation.

Every exponent in these definitions is a fixed nonnegative integer.
In particular, $q^6$ denotes an ordinary polynomial, evaluated by repeated
multiplication. There is no variable exponent, floor, factorial or
unbounded product in $P_{\Omega,r}$.

\section{A division-free center for the whole product}\label{sec:center}

The following short formula completely specifies the new shift. Put
\begin{align}
 e_2&=(kY)^2,\notag\\
 S_E&=(kY)\{\ell+k-(kY)(Y+1)\},\notag\\
 \Theta&=U^2b,\notag\\
 R_E&=3E+e_2,\notag\\
 H_E&=4E^2+3S_E^2+e_2^2,\notag\\
 \boxed{\displaystyle
 \Omega=U\left[2S_E(\Theta R_E-b+c)
           +A(4EV-\Theta H_E)\right].}\label{eq:Omega}
\end{align}
Finally, the requested equation is
\begin{equation}\label{eq:main}
 \boxed{\displaystyle
 P_{208}=Q\bigl(Z^2-16M^2GV^2\bigr)
               +\bigl(16M^2G-Z^2\bigr)=0.}
\end{equation}
Together with the finite coefficient recipe, these displayed definitions
are one complete polynomial specification in ten positive variables.
The supplied straight-line graph is a second, machine-readable
specification of exactly this polynomial.

\subsection{Where the formula comes from}

Introduce a formal indeterminate $t$, for the proof only, and write
\begin{align*}
 \mathcal E(t)&=e_2t^2+e_1t+e_0,\\
 e_1&=k^2Y(1-2Y-2Y^2),\\
 e_0&=1+k^2Y(Y+1)\{Y(Y+1)-1\},\\
 \mathcal V(t)&=\{1+2(t^2-1)U^2\mathcal E(t)\}b-c,\\
 \mathcal F(t)&=\mathcal E(t)\mathcal V(t).
\end{align*}
Then $\mathcal E(A)=E$, $\mathcal V(A)=V$ and $\mathcal F(A)=EV$.
When taking polynomial quotients, $Y,k,U,b,c$ are held as coefficient
parameters. Their eventual substitutions as polynomials in the original
ten variables do not invalidate a polynomial identity.

For $j=1,3$, let $Q_j$ be the value at $A$ of the ordinary polynomial
quotient of $\mathcal F(t)$ by $t^j$. Since $\mathcal E(t)(b-c)$ has
degree two and the other summand of $\mathcal F(t)$ is
$2U^2b(t^2-1)\mathcal E(t)^2$, every coefficient of $Q_3$ is even.
Thus the polynomial
\begin{equation}\label{eq:quotientcenter}
 U\left(4AEV-2Q_1-\frac{Q_3}{2}\right)
\end{equation}
has integer coefficients. Direct polynomial multiplication proves that
it is exactly~\eqref{eq:Omega}. The slash in~\eqref{eq:quotientcenter}
is explanatory notation only; equation~\eqref{eq:Omega} contains no
division operation.

The identity can also be checked without expanding a large expression.
Let $S=S_E$, $J_E=S(E+e_0)$ and
$D_E=e_2(2S-Ae_2)$. Since $E=AS+e_0$,
\[
 Q_1=S(b-c)+2\Theta(AE^2-J_E),\qquad
 \frac{Q_3}{2}=\Theta(J_E-D_E).
\]
Substitution into~\eqref{eq:quotientcenter} gives
\[
 \Omega=U\{4AEV-2S(b-c)-\Theta\mathcal H\},
 \quad
 \mathcal H=A(4E^2+3S^2+e_2^2)-2S(3E+e_2),
\]
and hence~\eqref{eq:Omega}.

The formal expansion
\[
 \sqrt{t^2-1}=t-\frac1{2t}-\frac1{8t^3}+O(t^{-5})
\]
suggests the center, but is not used as an analytic assumption in the
Diophantine proof. All required cancellations are exact polynomial
identities. The decisive change is to take the polynomial part of the
\emph{whole} expression $U\mathcal F(t)\sqrt{t^2-1}$, instead of first
approximating $g$ and then multiplying that approximation by $V$.

\section{Soundness before the signs of the encoded data are known}

This section states the inherited interfaces and proves the new signed
center estimate. The detailed parent arguments are retained in
\code{compact/parents/}, including its nested \code{proof_sources/}; their canonical
source paths and hashes are in the source ledger. The original unbounded
arguments are not replaced by finite numerical tests.

\subsection{Root extraction and the low-index exclusions}

At every positive input, independently of the equation,
\[
 q\ge32,\quad X,Y\ge q,\quad k\ge2,\quad c>kY,\quad
 U\ge c^2,\quad a>c,\quad f>2c,\quad M>0.
\]
Moreover,
\begin{equation}\label{eq:strongidentities}
 f^2=1+4aF_0,\qquad
 Q=\Delta(f^2-1)=4\Delta^2U^2G,\qquad
 f-1=2\Delta i^2c^4E.
\end{equation}
If $b=0$, then $V=-c$; otherwise $|V|\ge f-c>1$. At a zero of
\eqref{eq:main},
\[
 (Q-1)Z^2=16M^2G(QV^2-1).
\]
Since $Q$ is a multiple of $G$, both $Q-1$ and $QV^2-1$ are coprime
to $G$. Prime valuations show that $G$ is a square. The fixed factor
16 and the factor $M^2$ contribute only even valuations. Also
$\gcd(E,F_0)=1$, so
\[
 E=\tau^2,\quad F_0=f_0^2,\quad G=g^2,
 \qquad g=\tau f_0>0.
\]
The adjacent-coefficient Pell descent from the parent, applied initially
over the rationals, gives an integer $y>0$ such that
\begin{equation}\label{eq:adjacent}
 QV^2-(Q-1)y^2=1,\qquad Z/M=\varepsilon\,4gy,
 \quad \varepsilon\in\{-1,1\}.
\end{equation}
For completeness, if the rational ordinate has reduced denominator $m$,
then $m^2\mid Q-1$. The norm-preserving transformation
\[
 u'=(2Q-1)u-2(Q-1)y,\qquad
 y'=(2Q-1)y-2Qu
\]
strictly decreases $|u|>1$ after sign reflection, keeps its first
coordinate integral, and preserves the reduced denominator of the second.
It ends at $u=1,y=1$, proving $m=1$.

Put $I=2Ug$. The signed all-index folding theorem applies to
\[
 f^2-\Delta I^2=1,\quad Q=(\Delta I)^2,
 \quad c^2\mid I,\quad V\equiv-c\pmod f.
\]
It yields $c=\psi_A(p)$ for an odd positive $p$, where
$\psi_A(0)=0$, $\psi_A(1)=1$ and
$\psi_A(j+2)=2A\psi_A(j+1)-\psi_A(j)$.
The first norm yields $k=2\psi_{1+2XY^2}(m)$ for $m\ge1$.
The parent excludes $m=1$ and $p=3$, giving
\[
 c>8XY^3\ge8q^4,\quad c>X,\quad K>0,
 \quad c>4\Delta+3.
\]
These steps use no sign assumption about $B,C,R,b$.

The adjacent auxiliary index is odd. Index one would give $|V|=1$.
Index three would give one of the two impossible divisibilities
\[
 f(b-4\Delta f)=c-(4\Delta+3),\qquad
 f(b+4\Delta f)=c+(4\Delta+3),
\]
whose right sides lie strictly between zero and $f$. Thus its index is
at least five. Also $b=0$ is impossible, since then
$y^2-c^2=(c^2-1)/(Q-1)$ lies strictly between zero and one.
We obtain the strong absolute estimates
\begin{align}
 &\operatorname{sgn}(V)=\operatorname{sgn}(b),\quad
 |V|>8Q^2,\quad |b|>8f^3-\tfrac12>7f^3,\label{eq:bstrong}\\
 &g<f/6,\quad y<\tfrac32|V|,\quad
 |V|<\tfrac32f|b|,\quad AUE<\tfrac32g<f/4.
 \label{eq:absolutebounds}
\end{align}
The bound $g<f/6$ follows from
$a^2/g^2=\Delta U^2a/(1+a)>9$; the last bound uses
$G=E+\Delta(UE)^2$.

\subsection{A uniform polynomial-majorant lemma}

\begin{lemma}[The whole-product center is controlled]\label{lem:centerbound}
Under the preceding positive-input hypotheses and the consequences of
the Pell equations, before recovering the sign of $B$, the polynomial
$\Omega$ has the sign of $V$ and satisfies
\[
 3AUE|V|<|\Omega|<5AUE|V|.
\]
\end{lemma}
\begin{proof}
For a polynomial in the formal variable $t$, let a superscript $\#$ denote
the sum obtained by taking absolute values of its coefficients and then
evaluating at the positive real number $A$. This is proof notation, not
part of the Diophantine expression.

Put $\delta=(Y+1)/A$. Since $X,Y\ge32$, $\delta<1/16$. We have
\[
 |e_1|<2e_2(Y+1),\qquad 0<e_0\le e_2(Y+1)^2.
\]
Consequently
\[
 \mathcal E^{\#}(A)\le e_2(A+Y+1)^2,
 \qquad E\ge e_2A^2(1-2\delta),
\]
so $\mathcal E^{\#}(A)<2E$.

Using
$\mathcal F(t)=\mathcal E(t)(b-c)+2U^2b(t^2-1)\mathcal E(t)^2$
and $|b|\ge1$ gives
\[
 \mathcal F^{\#}(A)
 <2E(|b|+c)+8(A^2+1)U^2E^2|b|.
\]
Here $1+c<2U^2E$ and
$|V|>(f-c)|b|>\Delta U^2E|b|$. Therefore
\begin{equation}\label{eq:majorant}
 \frac{\mathcal F^{\#}(A)}{|EV|}
 <\frac4\Delta+\frac{8(A^2+1)}\Delta
 =8+\frac{20}{\Delta}<10.
\end{equation}
Deleting lower powers in a polynomial gives
$|Q_j|\le\mathcal F^{\#}(A)/A^j$. Equation~\eqref{eq:quotientcenter}
therefore implies
\[
 |\Omega-4AUEV|
 <10U|EV|\left(\frac2A+\frac1{2A^3}\right)
 <AU|EV|,
\]
because $A\ge16$ and $20/A^2+5/A^4<1$. This proves both the strict
sign assertion and the two bounds.
\end{proof}

\begin{corollary}[The square transport budget]\label{cor:budget}
Before assuming $B\ge0$,
\[
 4gy+|\Omega|<b^2/4.
\]
\end{corollary}
\begin{proof}
By~\eqref{eq:absolutebounds} and Lemma~\ref{lem:centerbound},
\[
 4gy+|\Omega|
 <6g|V|+5AUE|V|
 <\frac94f|V|
 <\frac{27}{8}f^2|b|.
\]
Since $|b|>7f^3$ and $f>2$, the ratio of the last bound to $b^2$
is less than $27/(56f)<1/4$.
\end{proof}

\subsection{Recovering the semantic signs without a circular argument}

For proof only, put
\[
 \rho=\frac{k-R-1-XY^2k(k^2-4)/6}{q^6},\quad
 \sigma=\frac{s}{q^2},\quad t_0=\frac{B}{h},\quad
 \kappa=\frac KH,\quad W=1+hb^2.
\]
The exact rational carrier identity is
\begin{equation}\label{eq:carrier}
 Z/M=n+\kappa+\sigma-Wt_0-\rho-\Omega.
\end{equation}
The inverse Lucas correction is crucial: $k\ge2$ implies
\[
 k-R-1-XY^2k(k^2-4)/6\le k-R-1<b.
\]
Thus $\rho<|b|/q^6<b^2/2$, uniformly for positive, zero or negative
$R$. If $B\le-1$, then $Wt_0<-b^2$. Since $\kappa,\sigma>0$,
\eqref{eq:adjacent}, \eqref{eq:carrier} and
Corollary~\ref{cor:budget} would give
\[
 n\le4gy+|\Omega|-\kappa-\sigma+Wt_0+\rho
 <\frac{b^2}{4}-b^2+\frac{b^2}{2}<0,
\]
contradicting $n>0$. Therefore $B\ge0$, and its defining equation
forces $C>0$.

Only now may the calibrated mask estimate $3q<R<q^4-2$ be used.
It gives $c>2R$ and hence $b,V>0$. The normalized-rank theorem and
the congruences
\[
 V\equiv-c\pmod f,\qquad V\equiv-R\pmod c
\]
identify $p=R$. The second congruence uses $f\equiv1\pmod c$ from
\eqref{eq:strongidentities}, not an assumed sign. The old signed cases
are thereby covered before the positive computation is reconstructed.

\section{Index recovery, completeness, and the logical endpoint}

\subsection{The inverse Lucas correction retains its exact denominator order}

Because $\Omega$ is an integer polynomial, $Z_0/M=\varepsilon4gy+\Omega$
is integral. Thus the parent inverse-correction arithmetic applies to
the unshifted carrier with exactly the same divisibility conditions.
For clarity, the essential order is summarized here.

The recurrence expansion, with $v=XY^2$, is
\[
 k=2\psi_{1+2v}(m)=2m+8v\binom{m+1}{3}+32v^2Z_m,
 \qquad Z_m\in\N.
\]
Writing $d=2m-R-1$ and
$e^*=k-R-1-vk(k^2-4)/6$, its exact cubic-difference identity gives
\[
 3(e^*-d)=4v^2\{24Z_m-A_mS_m\},
\]
where $A_m=\binom{m+1}{3}+4vZ_m$ and
$S_m=k^2+k(2m)+(2m)^2-4$. Since $v=q^3ws^2$, this proves
$q^4\mid(e^*-d)$ even at the prime 3. Integral $Z_0/M$ also gives
$q^4\mid e^*$. The first Pell base exceeds $A$, so $m<R$; the bound
$R<q^4$ now forces $d=0$.

One must not yet cancel the factor 3 at the full $q^6$ level.
First reduce the literal numerator modulo $H$ to obtain $H\mid6hK$.
The inherited uniform ratio bound and whole-product projection force
$X=2^R$, without canceling arbitrary factors of a congruence.
As $q\mid X$, $q$ is a power of two. Only then does the identity above
give $q^6\mid e^*$. It follows that $q^2\mid s$, $q^3\mid Y$,
$\kappa\in\Pos$ and $h\mid B$. These are exactly the native outer
conditions of the pinned calibration. This preserves the explicit
third-primary repair in source 164.

\subsection{Constructing positive witnesses from a refutation}

Conversely, a finite syntactic ZFC refutation supplies the parent's
canonical outer data with
\[
 X=2^R,
 \quad q^3\mid Y,
 \quad C>0,
 \quad t_0=B/h\in\N,
 \quad \rho,\sigma,\kappa\in\Z.
\]
The retained Pell-period construction gives positive $i,\eta,\zeta,T$
and positive $V,y$ satisfying the Pell equation in~\eqref{eq:adjacent}.
For the unsquared
trace, the old auxiliary coordinate changes by the explicitly integral
formula
\[
 T_{\rm new}=fT_{\rm old}-\frac{f-1}{c}R
            =\frac{V+c+Rf}{cf}>0.
\]
The quotient $(f-1)/c=2\Delta i^2c^3E$ is a polynomial.
The new center obeys $\Omega>0$. Define
\begin{equation}\label{eq:ncomplete}
 n=4gy+\Omega+(1+hb^2)t_0+\rho-\kappa-\sigma.
\end{equation}
Every term is an integer. The parent direct estimates give
$g>\kappa+\sigma+|\rho|$, including the potentially negative inverse
Lucas correction. Hence $n>0$, the carrier equals $4gy$, and the
literal polynomial vanishes. This constructs all ten positive witnesses;
it does not assert that a refutation or a zero actually exists.

\subsection{What the equivalence means internally}

The fixed tuple is the indexed native compilation of the calibrated
machine used by the pinned parents. Their outer reconstruction connects
a positive zero to a finite syntactic refutation. Composing the two
directions above gives, inside ZFC,
\[
 \exists\mathbf x\in\Pos^{10}\;P_{208}(\mathbf x)=0
 \quad\longleftrightarrow\quad
 \exists p\;\operatorname{Proof}_{\ZFC}(p,\bot).
\]
Negating proves Theorem~\ref{thm:main}. This reasoning uses finite coding,
integer arithmetic, recurrence identities and the specified inherited
proof transformations, all formalizable in ZFC. It is not an inference
from agreement only on standard numerals. Nor does it select coefficients
according to the unknown truth value of $\Con(\ZFC)$.

The direct quartics instead use the named interfaces
\code{KEndpoint.c222_iff_zf_refutation} and
\code{ConstructibleChoice.zf_refutation_iff_zfc_refutation}.
The new Diophantine transformations are ordinary mathematical proofs;
the inherited Lean/Rocq developments were not rebuilt in this session.

\section{Exact degree and the literal arithmetic count}\label{sec:degree}

The cancellation giving degree 208 must be checked on the entire polynomial.
Propagating degrees through the original squared expression alone would
give only an upper bound and miss the main result.

Assign the intermediate weights
\[
 \deg A=4,\quad\deg Y=2,\quad\deg k=1,
 \quad\deg U=7,\quad\deg b=4,\quad\deg c=3.
\]
An independent sparse expansion of the center gives
\begin{equation}\label{eq:defecttop}
 \topdeg{\Omega^2-16GV^2}{106}
 =8A^8Y^8k^8U^6b^2-48A^6Y^{12}k^8U^6b^2.
\end{equation}
The exact homogeneous component is meant here; these are not estimates
from a numerical ray. The full checked coefficient array of the defect
is included in the whole-product component.

The global polynomial identity
\begin{align}
 P_{208}={}&(Q-1)Z_0^2-2(Q-1)M\Omega Z_0\notag\\
 &+M^2\{Q(\Omega^2-16GV^2)-\Omega^2+16G\}
 \label{eq:wholeexpansion}
\end{align}
has successive upper degrees $122,177,208,152,72$.
In particular, using the square weight makes $\deg Z_0=21$;
retaining $bV$ here would make the cross term degree 213 and obscure the
lower-degree norm defect.

Set
\[
 q_0=\beta J,\qquad k_0=\eta+\zeta,\qquad
 B_0=k_0sT-q_0^2(q_0-F).
\]
Substitution of the actual leading components in~\eqref{eq:defecttop}
gives
\begin{equation}\label{eq:maintop}
 \boxed{\displaystyle
 \topdeg{P_{208}}{208}
 =9\,2^{11}q_0^{88}w^{18}s^{52}i^{10}k_0^{32}
 B_0^2(w^2-6s^2).}
\end{equation}
This is nonzero over the polynomial ring for the actual $\beta>0$.
Thus the exact degree is 208. The factors $B_0$ and $w^2-6s^2$ must
be retained; suppressing a tied leading term would give an incomplete
degree certificate.

The complete literal circuit has \PrimaryOps{} live binary operations:
\PrimaryMul{} multiplications, \PrimaryAdd{} additions and
\PrimarySub{} subtractions. Every gate, input and fixed port is used.
The independent checker binds its entire unchanged prefix to the frozen
85-operation parent, verifies the new center as a polynomial identity,
checks the altered carrier and finalizer, and establishes the exact
leading form. The reported count does not count a displayed abbreviation
as a free arithmetic operation.

\section{Direct quartics and nonnegativity on the zero set}\label{sec:quartics}

The direct component is logically independent of the native coefficient
tuple. Its source is package 164's C222 arithmetic construction, and its
projection technique descends from packages 159 and 163~\cite{p159,p163}. The fresh content
is the particular measured reductions and the natural-domain argument.
% BEGIN embedded direct_table.tex
% Generated by generate_inputs.py from the six packaged metrics files.
\begin{center}
\small
\begin{tabular}{llrrrr}
\toprule
Profile & Domain & Variables & Rows & Terms & Height\\
\midrule
Natural, principal & $\N$ & 203 & 155 & 8,237 & 799,349,232\\
Natural, lower height & $\N$ & 203 & 155 & 49,682 & 481,920\\
Natural, alternate basis & $\N$ & 203 & 155 & 8,263 & 207,116,893\\
Natural, intermediate & $\N$ & 204 & 156 & 8,213 & 130,686,374\\
Integer, principal & $\Z$ & 233 & 169 & 10,379 & 12,874\\
Integer, sparse & $\Z$ & 234 & 170 & 4,858 & 13,448\\
\bottomrule
\end{tabular}
\end{center}
All six polynomials have exact degree four. Terms and height refer
to the collected quartic; height is the largest absolute coefficient.

% END embedded direct_table.tex

The 203-natural equation trades coefficient height against the smaller
dimension; the row-basis variants make that tradeoff explicit. The
233-integer equation reduces the inspected 234-variable dimension record
in this direct family, but does not improve every coefficient measure.

\subsection{An exact quadratic projection}

Let $q_1,\ldots,q_m$ be integer quadratics. Suppose an integer vector $v$
is in the common kernel of their homogeneous quadratic Hessians, with
$v_p=1$. In coordinates $x=y+vz$, $y_p=0$, each row has the form
\[
 q_j(y+vz)=q_j(y)+a_jz,
 \qquad a_j\in\Z.
\]
If $a_r=\pm1$, eliminate $z=-a_rq_r(y)$. The remaining quadratics are
\[
 q_j(y)-a_ra_jq_r(y)=0\qquad(j\ne r).
\]
These maps are inverse on unrestricted integer zero sets. The independent
checker reconstructs the target by this row-combination identity; the
constructor uses shear and substitution. Natural domains require an
additional argument, which is not supplied by the Hessian calculation.

\subsection{Two canonical sign choices remove two natural coordinates}

In the pinned natural source, put
\[
 A_1=\code{mul.363},\quad B_1=\code{mul.369},
 \qquad Y_1=B_1-A_1.
\]
The row $A_1=(Sp+1)(1+\code{mul.361})$ restores $A_1\ge0$ and
$B_1=Y_1+A_1\ge0$ from every target natural zero. The source's canonical
CRT construction already chooses $B_1>A_1$, supplying the reverse
existential direction. The second pair is
\[
 C_1=\code{mul.372},\quad D_1=\code{mul.373},
 \qquad Y_2=C_1-D_1,
\]
with $D_1=(S+1)(1+\code{mul.361})$ and $C_1=Y_2+D_1$.
The same canonical representatives satisfy $C_1>D_1$.
These reductions give 204 nonnegative variables. They are not claimed to
be bijections on every arbitrary natural source zero.

\subsection{A retained equation certifies the third restoration}

Write $u=\code{mul.299}$, $v=\code{mul.304}$ and
$t=\code{add.302}$. Their defining rows are
$u=\Pi$ and $v=t(t+1)$, where
\begin{align*}
 \Pi={}&4556+4320r+1024r^2+512rs+128rt_1\\
       &{}+1080s+64s^2+32st_1+270t_1+4t_1^2,\\
 r&=\code{mul.232},\quad s=\code{mul.237},\quad
 t_1=\code{mul.243}.
\end{align*}
The common quadratic direction is $(u,v)\mapsto(u+z,v-4z)$.
Retain $w_1=v+4u$, eliminate $u$, and restore
\[
 u=\Pi,\qquad v=w_1-4\Pi.
\]
Although the latter polynomial has negative coefficients, its retained
row forces $w_1-4\Pi=t(t+1)\ge0$ on every target natural zero.
Conversely $w_1=v+4u\ge0$ on every old natural zero. Thus this third
step is a pointwise natural zero-set bijection and gives 203 variables.

The useful principle is that a restoration only has to be nonnegative
\emph{on the target zero set}. Requiring its coefficients all to be
nonnegative would miss this elimination. The package verifies the exact
retained row and the entire restoration, not just its expected shape.

\section{What was investigated and what is not claimed}

We investigated splitting the general C222 seed bundle into a shallow
guarded vacuity axiom and a fixed-label seed core. This initially looked
like a new semantic simplification. A source audit found exactly that
proposal already measured in package 156, with negative size results~\cite{p156}.
An updated trial still did not beat the current natural equation, so it
is not promoted as a new construction. This prevents an apparent change
of notation or an old negative experiment from being reported as a
breakthrough.

Whole-expression polynomial centers already occur in the
eleven-unrestricted-integer branch of packages 153, 154 and
165~\cite{signedcenters}. The new contribution here is the ten-positive-variable
specialization: the explicit quotient center for $EV$, its majorant bound
before recovery of the encoded signs, and the resulting literal degree-208
equation. Its unbounded proof is separate from the finite checks.
This endpoint and its specific argument appear new relative to the inspected
repository boundary; no exhaustive literature-priority claim is made.

There is no proof here that degree 208 is best for ten positive unknowns,
that 203 or 233 variables are best among direct quartics, or that any
one metric captures ``simplest.'' The small-variable forms retain
large fixed coefficients. The direct quartics have smaller coefficients
and many more variables. A future formalization must check the new
Diophantine bridge, not merely reuse the existing machine endpoint.

\section{Reproducibility and integration}

The ZIP contains this PDF and all native LaTeX sources; the ten-variable
graphs; four authenticated fixed numerals and their finite reconstruction;
the direct quadratic residual lists and collected quartic expansions;
ordinary proofs; independent coefficient/certificate checkers; source
ledgers, licenses, and recorded verification receipts.

The top-level command
\begin{quote}\code{python3 -B verify_all.py}\end{quote}
checks the new finite artifacts using only the Python standard library.
An optimized run is supported with \code{python3 -B -O verify_all.py}.
The report build command is \code{sh build_report.sh}; it requires an
ordinary TeX Live installation. Mathematical definitions are the
authenticated graph and coefficient arrays, not PDF metadata.

The final archive is an additive submission suitable for
\code{docs/incoming}. No remote repository content was modified.
The provenance boundary and internal AI-assisted reviews are disclosed;
they are not represented as outside human peer review.

\appendix
\section{The complete fixed coefficient recipe}\label{app:constants}

The following finite recipe is retained from the authenticated indexed
native source. It is a specification of four fixed integers, independent
of all quantified variables. Its bounded sums and fixed decoder do not
introduce variable exponents into the final polynomial. The executable
reconstructor verifies the exact raw bytes of all four integers.
% BEGIN embedded coefficient_recipe.tex
% Exact finite recipe extracted from the retained indexed_constants.tex.
% In this appendix K denotes the fixed port K_0, and all other symbols are constants.
\begingroup
\providecommand{\fl}[1]{\left\lfloor#1\right\rfloor}
\providecommand{\rg}[2]{\{#1,\ldots,#2\}}
In this appendix $K$ denotes the port $K_0$ in the new equations.
Every symbol introduced here is a fixed constant, independent of the ten
quantified variables. The formulas and 436-character seed are the retained
indexed BFS943 radix67 specification; they reconstruct the exact supplied
payloads, with $\beta=2^{590069}-1$.

\textbf{The four fixed integers.}
Every index, set, sum, product and floor below defines a constant, not
another unknown. Put
\[
\rho=2^{67},\quad B=\rho^{8807},\quad\beta=B-1,\quad
G_m(x)=\sum_{j=0}^{m-1}x^j.
\]
Define
\[
\begin{split}
\mathcal E={}&\rg0{942}\cup\rg{983}{992}\cup\rg{1003}{1011}\\
&\cup\rg{1013}{1031}\cup\{1063,1064,1095,2195,3296\}.
\end{split}
\]
Then
\begin{align*}
M_C={}&B-1-\sum_{e\in\mathcal E\setminus\{1\}}\rho^e-2\rho^{1063},\\
M_F={}&B-1+4+(2^{51}-2)(\rho^{4393}+\rho^{5495})\\
&+\rho^{1013}G_{10}(\rho)+\rho^{1065}G_9(\rho)
  +\rho^{2195}+\rho^{3296}\\
&+(2^{65}-1)\rho^{8792}G_{13}(\rho)+(2^{19}-1)\rho^{8805},\\
K={}&\rho^{30}+\rho^{62}+\rho^{1100}+\rho^{2201}+\rho^{8849}\\
&+(\rho^{4393}+\rho^{5495})\Phi(1/\rho),
\end{align*}
where
\[
\begin{split}
\Phi(x)={}&(1+2^{51})G_{943}(x)
 -2^{10}\sum_{i=0}^{942}v(t_i)x^i\\
&+\sum_{(b,e,m)\in\mathcal C}2^b x^e G_m(2x),\\
\mathcal C={}&\{(10,1023,8),(18,1013,9),(27,1003,8),\\
&(35,983,9),(44,1031,1),(45,1022,1),(46,1011,1),\\
&(47,992,1),(48,1095,1),(49,2195,1),(50,3296,1)\}.
\end{split}
\]
The rational intermediate in K cancels exactly; all four results are positive integers.

\textbf{The 943 indexed fixed tiles.}
Use the actions $u_j,v_j,d_j$ decoded below. Let
\[
\begin{gathered}
h(q,b)=2(q+1)+b,\quad
\mathcal R=\{u_j:d_j=1,\ 0\le j<430\},\\
\mathcal L=\{u_j:d_j=0,\ 0\le j<430\},\\
r_q=6+|\mathcal R\cap\rg0{q-1}|,\quad
l_q=205+|\mathcal L\cap\rg0{q-1}|.
\end{gathered}
\]
An interval ending at $-1$ is empty. The table defines $t_i$ for
$0\le i\le942$; all index ranges are fixed:
\[
\setlength{\arraycolsep}{1.4pt}\begin{array}{c|l}
\text{index }i&\text{tile }t_i=(w,n,e,s)\\\hline
0,1&(2,435,3,1),\ (5,0,0,0)\\
j+2\ (j=2q+b<430)&((1-d_j)l_{u_j},h(q,b),d_jr_{u_j},v_j)\\
432+b\ (b=0,1)&(0,h(215,b),0,h(215,b))\\
434+b\ (b=0,1)&(0,b,0,b)\\
424+2r_q+b\ (q\in\mathcal R,b=0,1)&(r_q,b,0,h(q,b))\\
424+2l_q+b\ (q\in\mathcal L,b=0,1)&(0,b,l_q,h(q,b))\\
924+j\ (0\le j<4)&(0,j+430\fl{j/2},0,438)\\
928&(2,434,2,0)\\
929&(3,436,4,4)\\
930&(4,436,4,0)\\
497+t\ (t=434,435,436)&(1,438,1,t)\\
\substack{934+3a+b-6\fl{a/4}\\a\in\{0,1,4\},\ b=0,1,2}&(a,437,b,437)
\end{array}
\]
For $t=(w,n,e,s)$ put
\[
v(t)=w+2^8n+2^{17}e+2^{25}s+127\,2^{34}.
\]

\textbf{A closed decoder for the 430 actions.}
Put $D_b(A,k)=\fl{A/b^k}\bmod b$, where
$a\bmod b=a-b\fl{a/b}$. For $0\le q<215$, $b=0,1$, $j=2q+b$:
\[
\begin{gathered}
N=\mathcal A\bmod2^{636},\quad
T_0=\fl{\mathcal A/2^{636}}\bmod2^{1675},\\
E_0=\fl{\mathcal A/2^{2311}},\\
\nu_j=D_2(N,j),\quad \sigma_j=\sum_{k=0}^{j-1}\nu_k,\\
m_q=D_2(N,430+q),\quad \tau_q=\sum_{k=0}^{q-1}m_k,\\
F_q=14(1-m_q)+m_qD_{16}(E_0,\tau_q),\\
u_j=\nu_j(1+\sigma_j)+(1-\nu_j)D_{216}(T_0,j-\sigma_j),\\
v_j=D_2(F_q,2b),\quad d_j=D_2(F_q,2b+1).
\end{gathered}
\]
Empty sums are zero. Here $(v_0,d_0,u_0)=(1,1,1)$; the constructor tile
uses this actual first target, including its state renaming.

\textbf{The integer $\mathcal A$.}
The alphabet, in order of digit values $0,\ldots,61$, is:\par
\noindent{\fontsize{7.7}{8.2}\selectfont\ttfamily 0123456789ABCDEFGHIJKLMNOPQRSTUVWXYZabcdefghijklmnopqrstuvwxyz}\par
Concatenate the following 436 characters as digits $e_0,\ldots,e_{435}$;
put $\mathcal A=\sum_{j=0}^{435}e_j62^{435-j}$.\par
\begingroup\fontsize{7.7}{8.2}\selectfont\ttfamily
BxEnhvy6NipzkWmROo0GsA3uK3twWLxscLq20UTumy5Cdu461LwD6g7RDNzF8Y\par
iYDFTpB8ehAypd7d6OsOL4APamnJeNXQypUezCmW9n4oXP6XfgmiWuAWZ8ZJ7U\par
OP7gu4kVdlGZqZHULQMar1dXNeJW4i7oNVI8usp8BlovjjZbwPCovZI328fRUq\par
4q92R81s5GjWQaQfGhkq9UYMK3rRMaBmsOIBkJqc2cj2VRIZ6mjzkBIOkW9SjZ\par
FzeA6J5u4j3sBL3SdgwXzM9mHvQivu8EJ1TBTpRhHDqZIbmYgwQ8VqTAFotl6F\par
We1YC554GKrVHsA8oUGKSA6F6k5KREBSD0Tgd7bFALY84LLbpw574BQyvdoHrl\par
0Oe35uwDf15DpT6VJWZBjSbHPpPAINYTte0eUxnTFTKBlGj8Am6sb1gLldh8WZ\par
8P\par
\endgroup

\endgroup

% END embedded coefficient_recipe.tex


\section{The principal literal arithmetic circuit}\label{app:circuit}

The following table prints the entire straight-line evaluator. A row
$z=x\mathbin\circ y$ introduces a polynomial abbreviation; only the ten
listed witnesses are quantified. The table is generated from, and checked
against, the canonical JSON graph.
% BEGIN embedded circuit.tex
% Generated by generate_inputs.py; all live source rows in order.
The last register, \code{N}, is $P_{208}$.
\begingroup
\fontsize{8.5}{10.2}\selectfont
\setlength{\tabcolsep}{6pt}
\setlength{\LTleft}{0pt}
\setlength{\LTright}{\fill}
\renewcommand{\arraystretch}{1.0}
\begin{longtable}{@{}rll@{}}
\toprule
Row & Register & Right-hand side\\
\midrule
\endfirsthead
\toprule
Row & Register & Right-hand side (continued)\\
\midrule
\endhead
\midrule
\multicolumn{3}{r}{Continued on next page}\\
\endfoot
\bottomrule
\endlastfoot
1 & \code{bJ} & \code{beta} \(\times\) \code{J}\\
2 & \code{q} & \code{bJ} \(+\) \(1\)\\
3 & \code{q2} & \code{q} \(\times\) \code{q}\\
4 & \code{q4} & \code{q2} \(\times\) \code{q2}\\
5 & \code{q6} & \code{q4} \(\times\) \code{q2}\\
6 & \code{sixq6} & \(6\) \(\times\) \code{q6}\\
7 & \code{Y} & \code{q} \(\times\) \code{s}\\
8 & \code{qmF} & \code{q} \(-\) \code{F}\\
9 & \code{C} & \code{qmF} \(-\) \code{alpha}\\
10 & \code{X} & \code{w} \(\times\) \code{q}\\
11 & \code{k} & \code{eta} \(+\) \code{zeta}\\
12 & \code{kY} & \code{k} \(\times\) \code{Y}\\
13 & \code{c} & \code{kY} \(+\) \code{eta}\\
14 & \code{XY} & \code{X} \(\times\) \code{Y}\\
15 & \code{Am1} & \code{XY} \(+\) \code{Y}\\
16 & \code{4Am1} & \(4\) \(\times\) \code{Am1}\\
17 & \code{Hp} & \code{4Am1} \(-\) \(1\)\\
18 & \code{Ap1} & \code{Am1} \(+\) \(2\)\\
19 & \code{V0k} & \code{XY} \(\times\) \code{kY}\\
20 & \code{V0kk} & \code{V0k} \(+\) \code{k}\\
21 & \code{V0k2} & \code{V0k} \(\times\) \code{V0kk}\\
22 & \code{Dtau} & \code{V0k2} \(+\) \(1\)\\
23 & \code{G1} & \code{bJ} \(\times\) \code{qmF}\\
24 & \code{G} & \code{G1} \(+\) \code{alpha}\\
25 & \code{q2m1} & \code{q2} \(-\) \(1\)\\
26 & \code{Gq} & \code{G} \(\times\) \code{q2m1}\\
27 & \code{qMF} & \code{q} \(\times\) \code{MF}\\
28 & \code{MCqMF} & \code{MC} \(+\) \code{qMF}\\
29 & \code{RJ} & \code{MCqMF} \(\times\) \code{J}\\
30 & \code{R} & \code{Gq} \(+\) \code{RJ}\\
31 & \code{Rp1} & \code{R} \(+\) \(1\)\\
32 & \code{e} & \code{k} \(-\) \code{Rp1}\\
33 & \code{ksq} & \code{k} \(\times\) \code{k}\\
34 & \code{ksqm4} & \code{ksq} \(-\) \(4\)\\
35 & \code{inverse_cubic} & \code{V0k} \(\times\) \code{ksqm4}\\
36 & \code{q4s} & \code{q4} \(\times\) \code{s}\\
37 & \code{e_sigma} & \code{e} \(-\) \code{q4s}\\
38 & \code{six_e_sigma} & \(6\) \(\times\) \code{e_sigma}\\
39 & \code{eclear} & \code{six_e_sigma} \(-\) \code{inverse_cubic}\\
40 & \code{K0w} & \code{K0} \(+\) \code{w}\\
41 & \code{K0wC} & \code{K0w} \(\times\) \code{C}\\
42 & \code{Bt} & \code{K0wC} \(-\) \code{F}\\
43 & \code{c2} & \code{c} \(\times\) \code{c}\\
44 & \code{U} & \code{i} \(\times\) \code{c2}\\
45 & \code{UE} & \code{U} \(\times\) \code{Dtau}\\
46 & \code{Delta} & \code{Am1} \(\times\) \code{Ap1}\\
47 & \code{SE} & \code{Delta} \(\times\) \code{UE}\\
48 & \code{DpI2} & \code{SE} \(\times\) \code{U}\\
49 & \code{Df} & \code{DpI2} \(+\) \(1\)\\
50 & \code{aF0} & \code{DpI2} \(\times\) \code{Df}\\
51 & \code{fouraF0} & \(4\) \(\times\) \code{aF0}\\
52 & \code{trace_f} & \code{DpI2} \(+\) \code{Df}\\
53 & \code{Q} & \code{Delta} \(\times\) \code{fouraF0}\\
54 & \code{cT} & \code{c} \(\times\) \code{T}\\
55 & \code{cTR} & \code{cT} \(-\) \code{R}\\
56 & \code{DfcTR} & \code{trace_f} \(\times\) \code{cTR}\\
57 & \code{V} & \code{DfcTR} \(-\) \code{c}\\
58 & \code{V2} & \code{V} \(\times\) \code{V}\\
59 & \code{P1mu2} & \code{bJ} \(\times\) \code{sixq6}\\
60 & \code{wp_A} & \code{Am1} \(+\) \(1\)\\
61 & \code{wp_a2} & \code{kY} \(\times\) \code{kY}\\
62 & \code{wp_Yp1} & \code{Y} \(+\) \(1\)\\
63 & \code{wp_kYYp1} & \code{kY} \(\times\) \code{wp_Yp1}\\
64 & \code{wp_sinner} & \code{V0kk} \(-\) \code{wp_kYYp1}\\
65 & \code{wp_S} & \code{kY} \(\times\) \code{wp_sinner}\\
66 & \code{wp_U2} & \code{U} \(\times\) \code{U}\\
67 & \code{wp_theta} & \code{wp_U2} \(\times\) \code{cTR}\\
68 & \code{wp_threeE} & \(3\) \(\times\) \code{Dtau}\\
69 & \code{wp_Zc} & \code{wp_threeE} \(+\) \code{wp_a2}\\
70 & \code{wp_thetaZc} & \code{wp_theta} \(\times\) \code{wp_Zc}\\
71 & \code{wp_bmc} & \code{cTR} \(-\) \code{c}\\
72 & \code{wp_Rc} & \code{wp_thetaZc} \(-\) \code{wp_bmc}\\
73 & \code{wp_E2} & \code{Dtau} \(\times\) \code{Dtau}\\
74 & \code{wp_fourE2} & \(4\) \(\times\) \code{wp_E2}\\
75 & \code{wp_S2} & \code{wp_S} \(\times\) \code{wp_S}\\
76 & \code{wp_threeS2} & \(3\) \(\times\) \code{wp_S2}\\
77 & \code{wp_a22} & \code{wp_a2} \(\times\) \code{wp_a2}\\
78 & \code{wp_sum1} & \code{wp_fourE2} \(+\) \code{wp_threeS2}\\
79 & \code{wp_Sc} & \code{wp_sum1} \(+\) \code{wp_a22}\\
80 & \code{wp_EV} & \code{Dtau} \(\times\) \code{V}\\
81 & \code{wp_fourEV} & \(4\) \(\times\) \code{wp_EV}\\
82 & \code{wp_thetaSc} & \code{wp_theta} \(\times\) \code{wp_Sc}\\
83 & \code{wp_Vdiff} & \code{wp_fourEV} \(-\) \code{wp_thetaSc}\\
84 & \code{wp_Apart} & \code{wp_A} \(\times\) \code{wp_Vdiff}\\
85 & \code{wp_SRc} & \code{wp_S} \(\times\) \code{wp_Rc}\\
86 & \code{wp_twoSRc} & \(2\) \(\times\) \code{wp_SRc}\\
87 & \code{wp_inner} & \code{wp_Apart} \(+\) \code{wp_twoSRc}\\
88 & \code{wp_W4} & \code{U} \(\times\) \code{wp_inner}\\
89 & \code{L} & \code{cTR} \(\times\) \code{cTR}\\
90 & \code{LBt} & \code{L} \(\times\) \code{Bt}\\
91 & \code{wp_raw_center} & \code{n} \(-\) \code{LBt}\\
92 & \code{center} & \code{wp_raw_center} \(-\) \code{wp_W4}\\
93 & \code{Hn} & \code{Hp} \(\times\) \code{center}\\
94 & \code{threec} & \(3\) \(\times\) \code{c}\\
95 & \code{twoX} & \(2\) \(\times\) \code{X}\\
96 & \code{gap} & \code{threec} \(-\) \code{twoX}\\
97 & \code{Kpos0} & \code{X} \(\times\) \code{gap}\\
98 & \code{Kpos} & \code{Kpos0} \(+\) \(2\)\\
99 & \code{lin} & \code{Hn} \(+\) \code{Kpos}\\
100 & \code{Zpart1} & \code{P1mu2} \(\times\) \code{lin}\\
101 & \code{mu2inner} & \code{sixq6} \(\times\) \code{Bt}\\
102 & \code{P1e} & \code{bJ} \(\times\) \code{eclear}\\
103 & \code{brk} & \code{mu2inner} \(+\) \code{P1e}\\
104 & \code{Zpart2} & \code{Hp} \(\times\) \code{brk}\\
105 & \code{Znum} & \code{Zpart1} \(-\) \code{Zpart2}\\
106 & \code{Mtot} & \code{P1mu2} \(\times\) \code{Hp}\\
107 & \code{M2} & \code{Mtot} \(\times\) \code{Mtot}\\
108 & \code{Dg} & \code{Dtau} \(\times\) \code{Df}\\
109 & \code{wp_G16} & \(16\) \(\times\) \code{Dg}\\
110 & \code{Z2} & \code{Znum} \(\times\) \code{Znum}\\
111 & \code{RG} & \code{M2} \(\times\) \code{wp_G16}\\
112 & \code{RGV2} & \code{RG} \(\times\) \code{V2}\\
113 & \code{inside} & \code{Z2} \(-\) \code{RGV2}\\
114 & \code{Qinside} & \code{Q} \(\times\) \code{inside}\\
115 & \code{correction} & \code{RG} \(-\) \code{Z2}\\
116 & \code{N} & \code{Qinside} \(+\) \code{correction}\\
\end{longtable}
\endgroup

% END embedded circuit.tex


\section{Literal residuals of the 203-natural quartic}\label{app:203}

Every coordinate below is independently in $\N$. The equation is exactly
$\sum_{j=0}^{154}q_j(x_0,\ldots,x_{202})^2=0$.
The complete expanded quartic in the archive is independently regenerated
from these residuals. An index denotes the canonical coordinate order in
\code{direct/natural203/equation.json}.
% BEGIN embedded residuals203.tex
% Generated by generate_inputs.py from direct/natural203/equation.json.
% The q_i order is exactly the residual-list order; indices are zero based.
\begingroup
\fontsize{9}{11}\selectfont
\allowdisplaybreaks[0]
\setlength{\abovedisplayskip}{3pt}
\setlength{\belowdisplayskip}{3pt}
\setlength{\abovedisplayshortskip}{3pt}
\setlength{\belowdisplayshortskip}{3pt}
\setlength{\jot}{1pt}
% residual 0; 3 literal terms
\begin{align*}
q_{0} &={}1 -x_{3}^{2} +x_{4}
\end{align*}
% residual 1; 6 literal terms
\begin{align*}
q_{1} &={}1 +x_{2} -x_{3} -x_{3}x_{5}\\
&\quad -x_{3}x_{6} -x_{5}
\end{align*}
% residual 2; 3 literal terms
\begin{align*}
q_{2} &={}-2\,x_{0}x_{2} -5\,x_{2} +x_{9}
\end{align*}
% residual 3; 5 literal terms
\begin{align*}
q_{3} &={}-1 +x_{1} -x_{2}x_{7} -x_{3}\\
&\quad -x_{7}
\end{align*}
% residual 4; 4 literal terms
\begin{align*}
q_{4} &={}-1 +x_{1} -x_{8} -x_{8}x_{9}
\end{align*}
% residual 5; 6 literal terms
\begin{align*}
q_{5} &={}1 -x_{4} -x_{5}x_{20} -x_{6}x_{20}\\
&\quad +x_{18} -x_{20}
\end{align*}
% residual 6; 8 literal terms
\begin{align*}
q_{6} &={}-5\,x_{10} -x_{10}^{2} -2\,x_{10}x_{11} +2\,x_{11}x_{12}\\
&\quad +5\,x_{12} +x_{12}^{2} +x_{21} -x_{25}
\end{align*}
% residual 7; 5 literal terms
\begin{align*}
q_{7} &={}-16 -8\,x_{11} -4\,x_{21} +x_{22}\\
&\quad -x_{35}
\end{align*}
% residual 8; 7 literal terms
\begin{align*}
q_{8} &={}-30 -11\,x_{10} -x_{10}^{2} -4\,x_{10}x_{22}\\
&\quad -22\,x_{22} -4\,x_{22}^{2} +x_{23}
\end{align*}
% residual 9; 7 literal terms
\begin{align*}
q_{9} &={}-380 -312\,x_{22} -64\,x_{22}^{2} -32\,x_{22}x_{23}\\
&\quad -78\,x_{23} -4\,x_{23}^{2} +x_{24}
\end{align*}
% residual 10; 7 literal terms
\begin{align*}
q_{10} &={}-6 -5\,x_{11} -x_{11}^{2} -2\,x_{11}x_{12}\\
&\quad -5\,x_{12} -x_{12}^{2} +x_{25}
\end{align*}
% residual 11; 9 literal terms
\begin{align*}
q_{11} &={}-6 +2\,x_{10}x_{11} -2\,x_{10}x_{12} -2\,x_{11}x_{12}\\
&\quad -10\,x_{12} -2\,x_{12}^{2} -x_{21} +x_{25}\\
&\quad +x_{26}
\end{align*}
% residual 12; 16 literal terms
\begin{align*}
q_{12} &={}-138 +8\,x_{10} -170\,x_{11} -30\,x_{11}^{2}\\
&\quad -48\,x_{11}x_{12} -32\,x_{11}x_{21} -24\,x_{11}x_{26} -48\,x_{12}\\
&\quad -32\,x_{12}x_{21} -100\,x_{21} -8\,x_{21}^{2} -16\,x_{21}x_{26}\\
&\quad +16\,x_{25} -40\,x_{26} +x_{27} -x_{31}
\end{align*}
% residual 13; 10 literal terms
\begin{align*}
q_{13} &={}-240 -32\,x_{10} -128\,x_{12} -64\,x_{25}\\
&\quad -6\,x_{27} -2\,x_{27}^{2} +x_{28} +6\,x_{31}\\
&\quad +2\,x_{31}^{2} -x_{32}
\end{align*}
% residual 14; 14 literal terms
\begin{align*}
q_{14} &={}-14 -4\,x_{10} -16\,x_{10}x_{25} -16\,x_{10}x_{26}\\
&\quad -4\,x_{11} -16\,x_{11}x_{25} -16\,x_{11}x_{26} -16\,x_{21}\\
&\quad -42\,x_{25} -4\,x_{25}^{2} -8\,x_{25}x_{26} -42\,x_{26}\\
&\quad -4\,x_{26}^{2} +x_{29}
\end{align*}
% residual 15; 16 literal terms
\begin{align*}
q_{15} &={}-1592 -2264\,x_{11} -800\,x_{11}^{2} -160\,x_{11}x_{21}\\
&\quad -640\,x_{11}x_{25} -160\,x_{11}x_{29} -220\,x_{21} -8\,x_{21}^{2}\\
&\quad -64\,x_{21}x_{25} -16\,x_{21}x_{29} -912\,x_{25} -128\,x_{25}^{2}\\
&\quad -64\,x_{25}x_{29} -228\,x_{29} -8\,x_{29}^{2} +x_{30}
\end{align*}
% residual 16; 16 literal terms
\begin{align*}
q_{16} &={}-180 +10\,x_{10} +6\,x_{10}^{2} -8\,x_{10}x_{25}\\
&\quad +10\,x_{11} +6\,x_{11}^{2} -8\,x_{11}x_{26} -192\,x_{12}\\
&\quad -48\,x_{12}^{2} -32\,x_{12}x_{25} -32\,x_{12}x_{26} -76\,x_{25}\\
&\quad -8\,x_{25}^{2} -76\,x_{26} -8\,x_{26}^{2} +x_{31}
\end{align*}
% residual 17; 13 literal terms
\begin{align*}
q_{17} &={}360 -40\,x_{10} -24\,x_{10}^{2} +32\,x_{10}x_{25}\\
&\quad +384\,x_{12} +96\,x_{12}^{2} +128\,x_{12}x_{25} +272\,x_{25}\\
&\quad +32\,x_{25}^{2} +32\,x_{26} -6\,x_{31} -2\,x_{31}^{2}\\
&\quad +x_{32}
\end{align*}
% residual 18; 11 literal terms
\begin{align*}
q_{18} &={}-110 -168\,x_{11} -64\,x_{11}^{2} -32\,x_{11}x_{21}\\
&\quad -32\,x_{11}x_{25} -42\,x_{21} -4\,x_{21}^{2} -8\,x_{21}x_{25}\\
&\quad -42\,x_{25} -4\,x_{25}^{2} +x_{33}
\end{align*}
% residual 19; 20 literal terms
\begin{align*}
q_{19} &={}-132 -220\,x_{10} -16\,x_{10}^{2} -128\,x_{10}x_{11}\\
&\quad -64\,x_{10}x_{21} -16\,x_{10}x_{26} -16\,x_{10}x_{33} -80\,x_{11}\\
&\quad -64\,x_{11}x_{26} -64\,x_{11}x_{33} -40\,x_{21} -32\,x_{21}x_{26}\\
&\quad -32\,x_{21}x_{33} -110\,x_{26} -4\,x_{26}^{2} -8\,x_{26}x_{33}\\
&\quad -110\,x_{33} -4\,x_{33}^{2} +x_{34} -4\,x_{39}
\end{align*}
% residual 20; 9 literal terms
\begin{align*}
q_{20} &={}-14 -8\,x_{11} +8\,x_{11}x_{12} +13\,x_{12}\\
&\quad +x_{12}^{2} +4\,x_{12}x_{21} -4\,x_{21} +x_{35}\\
&\quad -x_{61}
\end{align*}
% residual 21; 7 literal terms
\begin{align*}
q_{21} &={}-30 -11\,x_{10} -x_{10}^{2} -4\,x_{10}x_{35}\\
&\quad -22\,x_{35} -4\,x_{35}^{2} +x_{36}
\end{align*}
% residual 22; 7 literal terms
\begin{align*}
q_{22} &={}-380 -312\,x_{35} -64\,x_{35}^{2} -32\,x_{35}x_{36}\\
&\quad -78\,x_{36} -4\,x_{36}^{2} +x_{37}
\end{align*}
% residual 23; 7 literal terms
\begin{align*}
q_{23} &={}-30 -11\,x_{11} -x_{11}^{2} -4\,x_{11}x_{37}\\
&\quad -22\,x_{37} -4\,x_{37}^{2} +x_{38}
\end{align*}
% residual 24; 16 literal terms
\begin{align*}
q_{24} &={}-18 -8\,x_{10} -30\,x_{11} -34\,x_{11}^{2}\\
&\quad +48\,x_{11}x_{12} -32\,x_{11}x_{21} +24\,x_{11}x_{26} +48\,x_{12}\\
&\quad +32\,x_{12}x_{21} -8\,x_{21}^{2} +16\,x_{21}x_{26} -16\,x_{25}\\
&\quad +40\,x_{26} -x_{27} +x_{31} +x_{39}
\end{align*}
% residual 25; 15 literal terms
\begin{align*}
q_{25} &={}-782 -700\,x_{11} -144\,x_{11}^{2} -64\,x_{11}x_{21}\\
&\quad -80\,x_{11}x_{25} -80\,x_{11}x_{39} -200\,x_{21} -32\,x_{21}x_{25}\\
&\quad -32\,x_{21}x_{39} -150\,x_{25} -4\,x_{25}^{2} -8\,x_{25}x_{39}\\
&\quad -154\,x_{39} -4\,x_{39}^{2} +x_{40}
\end{align*}
% residual 26; 26 literal terms
\begin{align*}
q_{26} &={}-6272 -1020\,x_{10} -16\,x_{10}^{2} -512\,x_{10}x_{11}\\
&\quad -256\,x_{10}x_{21} -16\,x_{10}x_{26} -64\,x_{10}x_{39} -16\,x_{10}x_{40}\\
&\quad -3520\,x_{11} -256\,x_{11}x_{26} -1024\,x_{11}x_{39} -256\,x_{11}x_{40}\\
&\quad -1760\,x_{21} -128\,x_{21}x_{26} -512\,x_{21}x_{39} -128\,x_{21}x_{40}\\
&\quad -510\,x_{26} -4\,x_{26}^{2} -32\,x_{26}x_{39} -8\,x_{26}x_{40}\\
&\quad -2104\,x_{39} -64\,x_{39}^{2} -32\,x_{39}x_{40} -510\,x_{40}\\
&\quad -4\,x_{40}^{2} +x_{41}
\end{align*}
% residual 27; 7 literal terms
\begin{align*}
q_{27} &={}-6 -5\,x_{12} -x_{12}^{2} -2\,x_{12}x_{13}\\
&\quad -5\,x_{13} -x_{13}^{2} +x_{42}
\end{align*}
% residual 28; 20 literal terms
\begin{align*}
q_{28} &={}-224 -26\,x_{10} -2\,x_{10}^{2} -16\,x_{10}x_{13}\\
&\quad -8\,x_{10}x_{42} -288\,x_{12} -128\,x_{12}^{2} -64\,x_{12}x_{25}\\
&\quad -64\,x_{12}x_{26} -88\,x_{13} -32\,x_{13}^{2} -32\,x_{13}x_{42}\\
&\quad -76\,x_{25} -8\,x_{25}^{2} -16\,x_{25}x_{26} -68\,x_{26}\\
&\quad -8\,x_{26}^{2} -44\,x_{42} -8\,x_{42}^{2} +x_{43}
\end{align*}
% residual 29; 13 literal terms
\begin{align*}
q_{29} &={}256 +104\,x_{10} +8\,x_{10}^{2} +64\,x_{10}x_{13}\\
&\quad +32\,x_{10}x_{42} +352\,x_{13} +128\,x_{13}^{2} +128\,x_{13}x_{42}\\
&\quad +176\,x_{42} +32\,x_{42}^{2} -6\,x_{43} -2\,x_{43}^{2}\\
&\quad +x_{44}
\end{align*}
% residual 30; 7 literal terms
\begin{align*}
q_{30} &={}-42 -52\,x_{12} -16\,x_{12}^{2} -16\,x_{12}x_{25}\\
&\quad -26\,x_{25} -4\,x_{25}^{2} +x_{45}
\end{align*}
% residual 31; 7 literal terms
\begin{align*}
q_{31} &={}-30 -11\,x_{12} -x_{12}^{2} -4\,x_{12}x_{45}\\
&\quad -22\,x_{45} -4\,x_{45}^{2} +x_{46}
\end{align*}
% residual 32; 11 literal terms
\begin{align*}
q_{32} &={}-306 -35\,x_{11} -x_{11}^{2} -16\,x_{11}x_{45}\\
&\quad -4\,x_{11}x_{46} -280\,x_{45} -64\,x_{45}^{2} -32\,x_{45}x_{46}\\
&\quad -70\,x_{46} -4\,x_{46}^{2} +x_{47}
\end{align*}
% residual 33; 11 literal terms
\begin{align*}
q_{33} &={}-4556 -4320\,x_{45} -1024\,x_{45}^{2} -512\,x_{45}x_{46}\\
&\quad -128\,x_{45}x_{47} -1080\,x_{46} -64\,x_{46}^{2} -32\,x_{46}x_{47}\\
&\quad -270\,x_{47} -4\,x_{47}^{2} +x_{48}
\end{align*}
% residual 34; 7 literal terms
\begin{align*}
q_{34} &={}-30 -11\,x_{12} -x_{12}^{2} -4\,x_{12}x_{48}\\
&\quad -22\,x_{48} -4\,x_{48}^{2} +x_{49}
\end{align*}
% residual 35; 11 literal terms
\begin{align*}
q_{35} &={}-306 -35\,x_{11} -x_{11}^{2} -16\,x_{11}x_{48}\\
&\quad -4\,x_{11}x_{49} -280\,x_{48} -64\,x_{48}^{2} -32\,x_{48}x_{49}\\
&\quad -70\,x_{49} -4\,x_{49}^{2} +x_{50}
\end{align*}
% residual 36; 18 literal terms
\begin{align*}
q_{36} &={}-880 -1280\,x_{11} -480\,x_{11}^{2} -512\,x_{11}x_{25}\\
&\quad -128\,x_{11}x_{29} +74\,x_{12} +30\,x_{12}^{2} -32\,x_{12}x_{25}\\
&\quad -8\,x_{12}x_{29} -752\,x_{25} -128\,x_{25}^{2} -64\,x_{25}x_{29}\\
&\quad -180\,x_{29} -8\,x_{29}^{2} -32\,x_{48} -8\,x_{49}\\
&\quad -2\,x_{50} +x_{51}
\end{align*}
% residual 37; 7 literal terms
\begin{align*}
q_{37} &={}-264 -128\,x_{48} -32\,x_{49} -8\,x_{50}\\
&\quad -2\,x_{51} -2\,x_{51}^{2} +x_{52}
\end{align*}
% residual 38; 7 literal terms
\begin{align*}
q_{38} &={}-2 -3\,x_{10} -x_{10}^{2} -2\,x_{10}x_{52}\\
&\quad -3\,x_{52} -x_{52}^{2} +x_{53}
\end{align*}
% residual 39; 7 literal terms
\begin{align*}
q_{39} &={}-12 -28\,x_{52} -16\,x_{52}^{2} -16\,x_{52}x_{53}\\
&\quad -14\,x_{53} -4\,x_{53}^{2} +x_{54}
\end{align*}
% residual 40; 7 literal terms
\begin{align*}
q_{40} &={}-30 -11\,x_{12} -x_{12}^{2} -4\,x_{12}x_{54}\\
&\quad -22\,x_{54} -4\,x_{54}^{2} +x_{55}
\end{align*}
% residual 41; 11 literal terms
\begin{align*}
q_{41} &={}-306 -35\,x_{10} -x_{10}^{2} -16\,x_{10}x_{54}\\
&\quad -4\,x_{10}x_{55} -280\,x_{54} -64\,x_{54}^{2} -32\,x_{54}x_{55}\\
&\quad -70\,x_{55} -4\,x_{55}^{2} +x_{56}
\end{align*}
% residual 42; 11 literal terms
\begin{align*}
q_{42} &={}-4556 -4320\,x_{54} -1024\,x_{54}^{2} -512\,x_{54}x_{55}\\
&\quad -128\,x_{54}x_{56} -1080\,x_{55} -64\,x_{55}^{2} -32\,x_{55}x_{56}\\
&\quad -270\,x_{56} -4\,x_{56}^{2} +x_{57}
\end{align*}
% residual 43; 9 literal terms
\begin{align*}
q_{43} &={}-48 -24\,x_{12} -16\,x_{12}x_{57} -12\,x_{25}\\
&\quad -8\,x_{25}x_{57} -x_{45} -38\,x_{57} -4\,x_{57}^{2}\\
&\quad +x_{58}
\end{align*}
% residual 44; 20 literal terms
\begin{align*}
q_{44} &={}-9794 -70\,x_{11} -2\,x_{11}^{2} -32\,x_{11}x_{57}\\
&\quad -8\,x_{11}x_{58} -8640\,x_{48} -2048\,x_{48}^{2} -1024\,x_{48}x_{49}\\
&\quad -256\,x_{48}x_{50} -2160\,x_{49} -128\,x_{49}^{2} -64\,x_{49}x_{50}\\
&\quad -540\,x_{50} -8\,x_{50}^{2} -592\,x_{57} -128\,x_{57}^{2}\\
&\quad -64\,x_{57}x_{58} -148\,x_{58} -8\,x_{58}^{2} +x_{59}
\end{align*}
% residual 45; 13 literal terms
\begin{align*}
q_{45} &={}-18224 -17280\,x_{48} -4096\,x_{48}^{2} -2048\,x_{48}x_{49}\\
&\quad -512\,x_{48}x_{50} -4320\,x_{49} -256\,x_{49}^{2} -128\,x_{49}x_{50}\\
&\quad -1080\,x_{50} -16\,x_{50}^{2} -x_{59} -x_{59}^{2}\\
&\quad +x_{60}
\end{align*}
% residual 46; 12 literal terms
\begin{align*}
q_{46} &={}-30 -42\,x_{11} -14\,x_{11}^{2} -4\,x_{11}x_{12}\\
&\quad -16\,x_{11}x_{21} -3\,x_{12} +x_{12}^{2} -4\,x_{12}x_{21}\\
&\quad -26\,x_{21} -4\,x_{21}^{2} -2\,x_{25} +x_{61}
\end{align*}
% residual 47; 14 literal terms
\begin{align*}
q_{47} &={}522 +736\,x_{11} -160\,x_{11}x_{61} -70\,x_{12}\\
&\quad -2\,x_{12}^{2} -8\,x_{12}x_{60} +368\,x_{21} -80\,x_{21}x_{61}\\
&\quad -50\,x_{35} -148\,x_{60} -8\,x_{60}^{2} -228\,x_{61}\\
&\quad -8\,x_{61}^{2} +x_{62}
\end{align*}
% residual 48; 11 literal terms
\begin{align*}
q_{48} &={}-4800 -2944\,x_{11} +640\,x_{11}x_{61} -1472\,x_{21}\\
&\quad +320\,x_{21}x_{61} +200\,x_{35} +912\,x_{61} +32\,x_{61}^{2}\\
&\quad -6\,x_{62} -2\,x_{62}^{2} +x_{63}
\end{align*}
% residual 49; 4 literal terms
\begin{align*}
q_{49} &={}-2 -3\,x_{63} -x_{63}^{2} +x_{64}
\end{align*}
% residual 50; 3 literal terms
\begin{align*}
q_{50} &={}-x_{3} -x_{3}x_{15} +x_{65}
\end{align*}
% residual 51; 3 literal terms
\begin{align*}
q_{51} &={}-2\,x_{24}x_{65} -4\,x_{65} +x_{66}
\end{align*}
% residual 52; 2 literal terms
\begin{align*}
q_{52} &={}-x_{28}x_{65} +x_{67}
\end{align*}
% residual 53; 2 literal terms
\begin{align*}
q_{53} &={}-x_{30}x_{65} +x_{68}
\end{align*}
% residual 54; 2 literal terms
\begin{align*}
q_{54} &={}-x_{32}x_{65} +x_{69}
\end{align*}
% residual 55; 6 literal terms
\begin{align*}
q_{55} &={}-64\,x_{11}x_{65} -32\,x_{21}x_{65} -8\,x_{33}x_{65} -2\,x_{34}x_{65}\\
&\quad -80\,x_{65} +x_{70}
\end{align*}
% residual 56; 4 literal terms
\begin{align*}
q_{56} &={}-4\,x_{11}x_{65} -2\,x_{38}x_{65} -4\,x_{65} +x_{71}
\end{align*}
% residual 57; 7 literal terms
\begin{align*}
q_{57} &={}-256\,x_{11}x_{65} -128\,x_{21}x_{65} -32\,x_{39}x_{65} -8\,x_{40}x_{65}\\
&\quad -2\,x_{41}x_{65} -480\,x_{65} +x_{72}
\end{align*}
% residual 58; 2 literal terms
\begin{align*}
q_{58} &={}-x_{44}x_{65} +x_{73}
\end{align*}
% residual 59; 3 literal terms
\begin{align*}
q_{59} &={}-2\,x_{64}x_{65} -4\,x_{65} +x_{74}
\end{align*}
% residual 60; 3 literal terms
\begin{align*}
q_{60} &={}-x_{3} -x_{3}x_{11} +x_{75}
\end{align*}
% residual 61; 3 literal terms
\begin{align*}
q_{61} &={}-x_{3} -x_{3}x_{10} +x_{76}
\end{align*}
% residual 62; 3 literal terms
\begin{align*}
q_{62} &={}2\,x_{3} +x_{77} -x_{80}
\end{align*}
% residual 63; 5 literal terms
\begin{align*}
q_{63} &={}-1 -x_{14} -x_{14}x_{76} -x_{76}\\
&\quad +x_{78}
\end{align*}
% residual 64; 7 literal terms
\begin{align*}
q_{64} &={}1 +x_{15} +x_{15}x_{75} +x_{75}\\
&\quad -x_{77}x_{78} -x_{78} +x_{79}
\end{align*}
% residual 65; 4 literal terms
\begin{align*}
q_{65} &={}-6\,x_{3} -4\,x_{3}x_{11} -2\,x_{3}x_{21} +x_{80}
\end{align*}
% residual 66; 7 literal terms
\begin{align*}
q_{66} &={}x_{14} +x_{14}x_{75} -x_{15} -x_{15}x_{80}\\
&\quad +x_{75} -x_{80} +x_{81}
\end{align*}
% residual 67; 3 literal terms
\begin{align*}
q_{67} &={}-2\,x_{2} -x_{2}x_{18} +x_{82}
\end{align*}
% residual 68; 5 literal terms
\begin{align*}
q_{68} &={}-1 -x_{16} -x_{16}x_{82} -x_{82}\\
&\quad +x_{83}
\end{align*}
% residual 69; 7 literal terms
\begin{align*}
q_{69} &={}-1 -x_{2} -x_{2}x_{17} -x_{17}\\
&\quad -x_{17}x_{82} -x_{82} +x_{84}
\end{align*}
% residual 70; 20 literal terms
\begin{align*}
q_{70} &={}2 -4\,x_{1} +2\,x_{1}^{2} -2\,x_{1}x_{14}\\
&\quad -2\,x_{1}x_{15} -2\,x_{1}x_{83} -2\,x_{1}x_{84} +2\,x_{14}\\
&\quad +x_{14}^{2} +2\,x_{14}x_{83} +2\,x_{15} +x_{15}^{2}\\
&\quad +2\,x_{15}x_{84} -x_{19}x_{85} +2\,x_{83} +x_{83}^{2}\\
&\quad +2\,x_{84} +x_{84}^{2} +x_{201}^{2} +x_{202}^{2}
\end{align*}
% residual 71; 11 literal terms
\begin{align*}
q_{71} &={}-x_{10}x_{19} -x_{11}x_{19} -x_{12}x_{19} -x_{13}x_{19}\\
&\quad -x_{14}x_{19} -x_{15}x_{19} -x_{16}x_{19} -x_{17}x_{19}\\
&\quad -x_{18}x_{19} -10\,x_{19} +x_{86}
\end{align*}
% residual 72; 11 literal terms
\begin{align*}
q_{72} &={}1 -x_{87} -x_{87}x_{132} -x_{87}x_{133}\\
&\quad -x_{87}x_{183} +x_{138} +x_{139} +x_{145}\\
&\quad +x_{150} +x_{155} +x_{186}
\end{align*}
% residual 73; 5 literal terms
\begin{align*}
q_{73} &={}-x_{19} -x_{19}x_{172} -x_{19}x_{173} -x_{19}x_{178}\\
&\quad +x_{88}
\end{align*}
% residual 74; 72 literal terms
\begin{align*}
q_{74} &={}-34 -2\,x_{0} -2\,x_{0}x_{91} -x_{1}\\
&\quad -x_{1}x_{91} -2\,x_{2} -2\,x_{2}x_{91} -x_{4}\\
&\quad -x_{4}x_{91} -x_{5} -x_{5}x_{91} -x_{6}\\
&\quad -x_{6}x_{91} -x_{10} -2\,x_{10}x_{91} -x_{11}\\
&\quad -334\,x_{11}x_{91} -x_{12} -x_{12}x_{91} -x_{13}\\
&\quad -x_{13}x_{91} -x_{14} -x_{14}x_{91} -x_{15}\\
&\quad -x_{15}x_{91} -x_{16} -x_{16}x_{91} -x_{17}\\
&\quad -x_{17}x_{91} -x_{18} -x_{18}x_{91} -x_{19}\\
&\quad -x_{19}x_{91} -164\,x_{21}x_{91} -2\,x_{24}x_{91} -x_{28}x_{91}\\
&\quad -x_{30}x_{91} -x_{32}x_{91} -8\,x_{33}x_{91} -2\,x_{34}x_{91}\\
&\quad -2\,x_{38}x_{91} -32\,x_{39}x_{91} -8\,x_{40}x_{91} -2\,x_{41}x_{91}\\
&\quad -x_{44}x_{91} -2\,x_{64}x_{91} -x_{86} -x_{86}x_{91}\\
&\quad -x_{87} -x_{87}x_{91} -619\,x_{91} -x_{91}x_{132}\\
&\quad -x_{91}x_{133} -x_{91}x_{160} -x_{91}x_{161} -x_{91}x_{167}\\
&\quad -x_{91}x_{172} -x_{91}x_{173} -x_{91}x_{178} -x_{91}x_{183}\\
&\quad +x_{92} +4\,x_{96}x_{107} -x_{98} -x_{132}\\
&\quad -x_{133} -x_{160} -x_{161} -x_{167}\\
&\quad -x_{172} -x_{173} -x_{178} -x_{183}
\end{align*}
% residual 75; 42 literal terms
\begin{align*}
q_{75} &={}-2\,x_{0}x_{92} -x_{1}x_{92} -2\,x_{2}x_{92} -x_{4}x_{92}\\
&\quad -x_{5}x_{92} -x_{6}x_{92} -2\,x_{10}x_{92} -334\,x_{11}x_{92}\\
&\quad -x_{12}x_{92} -x_{13}x_{92} -x_{14}x_{92} -x_{15}x_{92}\\
&\quad -x_{16}x_{92} -x_{17}x_{92} -x_{18}x_{92} -x_{19}x_{92}\\
&\quad -164\,x_{21}x_{92} -2\,x_{24}x_{92} -x_{28}x_{92} -x_{30}x_{92}\\
&\quad -x_{32}x_{92} -8\,x_{33}x_{92} -2\,x_{34}x_{92} -2\,x_{38}x_{92}\\
&\quad -32\,x_{39}x_{92} -8\,x_{40}x_{92} -2\,x_{41}x_{92} -x_{44}x_{92}\\
&\quad -2\,x_{64}x_{92} -x_{86}x_{92} -x_{87}x_{92} -620\,x_{92}\\
&\quad -x_{92}x_{132} -x_{92}x_{133} -x_{92}x_{160} -x_{92}x_{161}\\
&\quad -x_{92}x_{167} -x_{92}x_{172} -x_{92}x_{173} -x_{92}x_{178}\\
&\quad -x_{92}x_{183} +x_{93}
\end{align*}
% residual 76; 4 literal terms
\begin{align*}
q_{76} &={}-1 +x_{89}^{2} +x_{92}^{2} -x_{93}^{2}
\end{align*}
% residual 77; 6 literal terms
\begin{align*}
q_{77} &={}-1 -2\,x_{87} -x_{88}x_{94} +2\,x_{89}\\
&\quad +2\,x_{90} +x_{121}
\end{align*}
% residual 78; 2 literal terms
\begin{align*}
q_{78} &={}-x_{97}^{2} +x_{188}
\end{align*}
% residual 79; 8 literal terms
\begin{align*}
q_{79} &={}1 -x_{89} +x_{89}x_{188} -x_{90}\\
&\quad +x_{90}x_{188} -4\,x_{96}x_{188} +x_{99} -x_{188}
\end{align*}
% residual 80; 3 literal terms
\begin{align*}
q_{80} &={}-x_{89}x_{96} -x_{90}x_{96} +x_{100}
\end{align*}
% residual 81; 4 literal terms
\begin{align*}
q_{81} &={}-1 +x_{95}^{2} +x_{96}^{2} -x_{100}^{2}
\end{align*}
% residual 82; 3 literal terms
\begin{align*}
q_{82} &={}-2\,x_{89}x_{105} -2\,x_{90}x_{105} +x_{101}
\end{align*}
% residual 83; 4 literal terms
\begin{align*}
q_{83} &={}-1 +x_{97}^{2} -x_{101}^{2} +4\,x_{105}^{2}
\end{align*}
% residual 84; 3 literal terms
\begin{align*}
q_{84} &={}-x_{98} -x_{98}x_{99} +x_{102}
\end{align*}
% residual 85; 6 literal terms
\begin{align*}
q_{85} &={}-1 +x_{95}^{2} +2\,x_{95}x_{108} +x_{98}^{2}\\
&\quad -x_{102}^{2} +x_{108}^{2}
\end{align*}
% residual 86; 2 literal terms
\begin{align*}
q_{86} &={}-x_{96}^{2} +x_{104}
\end{align*}
% residual 87; 3 literal terms
\begin{align*}
q_{87} &={}-x_{103}x_{104} -x_{104} +x_{105}
\end{align*}
% residual 88; 2 literal terms
\begin{align*}
q_{88} &={}-x_{97}x_{106} +x_{108}
\end{align*}
% residual 89; 18 literal terms
\begin{align*}
q_{89} &={}-585 -x_{10} -333\,x_{11} -164\,x_{21}\\
&\quad -2\,x_{24} -x_{28} -x_{30} -x_{32}\\
&\quad -8\,x_{33} -2\,x_{34} -2\,x_{38} -32\,x_{39}\\
&\quad -8\,x_{40} -2\,x_{41} -x_{44} -2\,x_{64}\\
&\quad -4\,x_{96}x_{107} +x_{98}
\end{align*}
% residual 90; 2 literal terms
\begin{align*}
q_{90} &={}-x_{109}x_{110} +x_{111}
\end{align*}
% residual 91; 5 literal terms
\begin{align*}
q_{91} &={}-1 -2\,x_{89}x_{111} -2\,x_{90}x_{111} +x_{109}^{2}\\
&\quad +x_{110}^{2}
\end{align*}
% residual 92; 5 literal terms
\begin{align*}
q_{92} &={}-1 -x_{4}^{2} +2\,x_{4}x_{89} +2\,x_{4}x_{90}\\
&\quad -x_{112}
\end{align*}
% residual 93; 5 literal terms
\begin{align*}
q_{93} &={}1 +x_{2} -x_{4}x_{109} +x_{110}\\
&\quad +x_{112}x_{113}
\end{align*}
% residual 94; 2 literal terms
\begin{align*}
q_{94} &={}-x_{114}x_{115} +x_{116}
\end{align*}
% residual 95; 5 literal terms
\begin{align*}
q_{95} &={}-1 -2\,x_{89}x_{116} -2\,x_{90}x_{116} +x_{114}^{2}\\
&\quad +x_{115}^{2}
\end{align*}
% residual 96; 13 literal terms
\begin{align*}
q_{96} &={}-2 -2\,x_{5} -x_{5}^{2} -2\,x_{5}x_{6}\\
&\quad +2\,x_{5}x_{89} +2\,x_{5}x_{90} -2\,x_{6} -x_{6}^{2}\\
&\quad +2\,x_{6}x_{89} +2\,x_{6}x_{90} +2\,x_{89} +2\,x_{90}\\
&\quad -x_{117}
\end{align*}
% residual 97; 6 literal terms
\begin{align*}
q_{97} &={}-x_{5}x_{114} -x_{6}x_{114} +x_{19} -x_{114}\\
&\quad +x_{115} +x_{117}x_{118}
\end{align*}
% residual 98; 5 literal terms
\begin{align*}
q_{98} &={}-1 -x_{86}^{2} +2\,x_{86}x_{89} +2\,x_{86}x_{90}\\
&\quad -x_{119}
\end{align*}
% residual 99; 4 literal terms
\begin{align*}
q_{99} &={}x_{86}x_{114} -x_{87} -x_{115} -x_{119}x_{120}
\end{align*}
% residual 100; 5 literal terms
\begin{align*}
q_{100} &={}1 +x_{87}^{2} -2\,x_{87}x_{89} -2\,x_{87}x_{90}\\
&\quad +x_{121}
\end{align*}
% residual 101; 7 literal terms
\begin{align*}
q_{101} &={}-1 +x_{87}x_{109} -x_{110} -x_{121}x_{122}\\
&\quad -x_{132} -x_{133} -x_{183}
\end{align*}
% residual 102; 2 literal terms
\begin{align*}
q_{102} &={}-x_{123}x_{124} +x_{125}
\end{align*}
% residual 103; 5 literal terms
\begin{align*}
q_{103} &={}-1 -2\,x_{89}x_{125} -2\,x_{90}x_{125} +x_{123}^{2}\\
&\quad +x_{124}^{2}
\end{align*}
% residual 104; 7 literal terms
\begin{align*}
q_{104} &={}-1 +x_{87}x_{123} -x_{121}x_{126} -x_{124}\\
&\quad -x_{160} -x_{161} -x_{167}
\end{align*}
% residual 105; 7 literal terms
\begin{align*}
q_{105} &={}1 -x_{87}x_{114} +x_{115} +x_{121}x_{127}\\
&\quad +x_{172} +x_{173} +x_{178}
\end{align*}
% residual 106; 2 literal terms
\begin{align*}
q_{106} &={}-x_{128}x_{129} +x_{130}
\end{align*}
% residual 107; 5 literal terms
\begin{align*}
q_{107} &={}-1 -2\,x_{89}x_{130} -2\,x_{90}x_{130} +x_{128}^{2}\\
&\quad +x_{129}^{2}
\end{align*}
% residual 108; 22 literal terms
\begin{align*}
q_{108} &={}1 +x_{5} +x_{5}x_{132} +x_{5}x_{133}\\
&\quad +x_{5}x_{183} +x_{6} +x_{6}x_{132} +x_{6}x_{133}\\
&\quad +x_{6}x_{183} -x_{87}x_{128} -x_{128} +x_{129}\\
&\quad +x_{131} +x_{131}x_{138} +x_{131}x_{139} +x_{131}x_{145}\\
&\quad +x_{131}x_{150} +x_{131}x_{155} +x_{131}x_{186} +2\,x_{132}\\
&\quad +x_{133} +x_{183}
\end{align*}
% residual 109; 2 literal terms
\begin{align*}
q_{109} &={}-x_{134}x_{135} +x_{136}
\end{align*}
% residual 110; 5 literal terms
\begin{align*}
q_{110} &={}-1 -2\,x_{89}x_{136} -2\,x_{90}x_{136} +x_{134}^{2}\\
&\quad +x_{135}^{2}
\end{align*}
% residual 111; 8 literal terms
\begin{align*}
q_{111} &={}-x_{19} -x_{19}x_{138} -x_{19}x_{139} -x_{19}x_{145}\\
&\quad -x_{19}x_{150} -x_{19}x_{155} -x_{19}x_{186} +x_{140}
\end{align*}
% residual 112; 5 literal terms
\begin{align*}
q_{112} &={}x_{87}x_{134} +x_{134} -x_{135} -x_{137}x_{140}\\
&\quad -x_{138}
\end{align*}
% residual 113; 2 literal terms
\begin{align*}
q_{113} &={}-x_{141}x_{142} +x_{143}
\end{align*}
% residual 114; 5 literal terms
\begin{align*}
q_{114} &={}-1 -2\,x_{89}x_{143} -2\,x_{90}x_{143} +x_{141}^{2}\\
&\quad +x_{142}^{2}
\end{align*}
% residual 115; 5 literal terms
\begin{align*}
q_{115} &={}-x_{87}x_{141} +x_{140}x_{144} -x_{141} +x_{142}\\
&\quad +x_{145}
\end{align*}
% residual 116; 2 literal terms
\begin{align*}
q_{116} &={}-x_{146}x_{147} +x_{148}
\end{align*}
% residual 117; 5 literal terms
\begin{align*}
q_{117} &={}-1 -2\,x_{89}x_{148} -2\,x_{90}x_{148} +x_{146}^{2}\\
&\quad +x_{147}^{2}
\end{align*}
% residual 118; 5 literal terms
\begin{align*}
q_{118} &={}-x_{87}x_{146} +x_{140}x_{149} -x_{146} +x_{147}\\
&\quad +x_{150}
\end{align*}
% residual 119; 2 literal terms
\begin{align*}
q_{119} &={}-x_{151}x_{152} +x_{153}
\end{align*}
% residual 120; 5 literal terms
\begin{align*}
q_{120} &={}-1 -2\,x_{89}x_{153} -2\,x_{90}x_{153} +x_{151}^{2}\\
&\quad +x_{152}^{2}
\end{align*}
% residual 121; 5 literal terms
\begin{align*}
q_{121} &={}-x_{87}x_{151} +x_{140}x_{154} -x_{151} +x_{152}\\
&\quad +x_{155}
\end{align*}
% residual 122; 2 literal terms
\begin{align*}
q_{122} &={}-x_{156}x_{157} +x_{158}
\end{align*}
% residual 123; 5 literal terms
\begin{align*}
q_{123} &={}-1 -2\,x_{89}x_{158} -2\,x_{90}x_{158} +x_{156}^{2}\\
&\quad +x_{157}^{2}
\end{align*}
% residual 124; 5 literal terms
\begin{align*}
q_{124} &={}-x_{19} -x_{19}x_{160} -x_{19}x_{161} -x_{19}x_{167}\\
&\quad +x_{162}
\end{align*}
% residual 125; 5 literal terms
\begin{align*}
q_{125} &={}x_{87}x_{156} +x_{156} -x_{157} -x_{159}x_{162}\\
&\quad -x_{160}
\end{align*}
% residual 126; 2 literal terms
\begin{align*}
q_{126} &={}-x_{163}x_{164} +x_{165}
\end{align*}
% residual 127; 5 literal terms
\begin{align*}
q_{127} &={}-1 -2\,x_{89}x_{165} -2\,x_{90}x_{165} +x_{163}^{2}\\
&\quad +x_{164}^{2}
\end{align*}
% residual 128; 5 literal terms
\begin{align*}
q_{128} &={}-x_{87}x_{163} +x_{162}x_{166} -x_{163} +x_{164}\\
&\quad +x_{167}
\end{align*}
% residual 129; 2 literal terms
\begin{align*}
q_{129} &={}-x_{168}x_{169} +x_{170}
\end{align*}
% residual 130; 5 literal terms
\begin{align*}
q_{130} &={}-1 -2\,x_{89}x_{170} -2\,x_{90}x_{170} +x_{168}^{2}\\
&\quad +x_{169}^{2}
\end{align*}
% residual 131; 5 literal terms
\begin{align*}
q_{131} &={}-x_{87}x_{168} +x_{88}x_{171} -x_{168} +x_{169}\\
&\quad +x_{172}
\end{align*}
% residual 132; 2 literal terms
\begin{align*}
q_{132} &={}-x_{174}x_{175} +x_{176}
\end{align*}
% residual 133; 5 literal terms
\begin{align*}
q_{133} &={}-1 -2\,x_{89}x_{176} -2\,x_{90}x_{176} +x_{174}^{2}\\
&\quad +x_{175}^{2}
\end{align*}
% residual 134; 5 literal terms
\begin{align*}
q_{134} &={}x_{87}x_{174} -x_{88}x_{177} +x_{174} -x_{175}\\
&\quad -x_{178}
\end{align*}
% residual 135; 2 literal terms
\begin{align*}
q_{135} &={}-x_{179}x_{180} +x_{181}
\end{align*}
% residual 136; 5 literal terms
\begin{align*}
q_{136} &={}-1 -2\,x_{89}x_{181} -2\,x_{90}x_{181} +x_{179}^{2}\\
&\quad +x_{180}^{2}
\end{align*}
% residual 137; 5 literal terms
\begin{align*}
q_{137} &={}-x_{19} -x_{19}x_{132} -x_{19}x_{133} -x_{19}x_{183}\\
&\quad +x_{184}
\end{align*}
% residual 138; 5 literal terms
\begin{align*}
q_{138} &={}x_{87}x_{179} +x_{179} -x_{180} -x_{182}x_{184}\\
&\quad -x_{183}
\end{align*}
% residual 139; 6 literal terms
\begin{align*}
q_{139} &={}x_{87}x_{96} +x_{95} +x_{96} -x_{100}\\
&\quad -x_{140}x_{185} -x_{186}
\end{align*}
% residual 140; 68 literal terms
\begin{align*}
q_{140} &={}-26 -16\,x_{0} -4\,x_{0}^{2} +4\,x_{0}x_{109}\\
&\quad -x_{1}^{2} +2\,x_{1}x_{114} -2\,x_{2} -x_{2}^{2}\\
&\quad +2\,x_{2}x_{123} -x_{4}^{2} +2\,x_{4}x_{128} -2\,x_{10}\\
&\quad -x_{10}^{2} +2\,x_{10}x_{134} -2\,x_{11} -x_{11}^{2}\\
&\quad +2\,x_{11}x_{141} -2\,x_{12} -x_{12}^{2} +2\,x_{12}x_{146}\\
&\quad -2\,x_{13} -x_{13}^{2} +2\,x_{13}x_{151} -2\,x_{14}\\
&\quad -x_{14}^{2} +2\,x_{14}x_{156} -2\,x_{15} -x_{15}^{2}\\
&\quad +2\,x_{15}x_{163} -2\,x_{16} -x_{16}^{2} +2\,x_{16}x_{168}\\
&\quad -2\,x_{17} -x_{17}^{2} +2\,x_{17}x_{174} -2\,x_{18}\\
&\quad -x_{18}^{2} +2\,x_{18}x_{179} -x_{89}x_{187} -x_{90}x_{187}\\
&\quad +x_{96} +4\,x_{96}x_{107} -x_{98} +8\,x_{109}\\
&\quad -x_{109}^{2} -x_{114}^{2} +2\,x_{123} -x_{123}^{2}\\
&\quad -x_{128}^{2} +2\,x_{134} -x_{134}^{2} +2\,x_{141}\\
&\quad -x_{141}^{2} +2\,x_{146} -x_{146}^{2} +2\,x_{151}\\
&\quad -x_{151}^{2} +2\,x_{156} -x_{156}^{2} +2\,x_{163}\\
&\quad -x_{163}^{2} +2\,x_{168} -x_{168}^{2} +2\,x_{174}\\
&\quad -x_{174}^{2} +2\,x_{179} -x_{179}^{2} +x_{187}
\end{align*}
% residual 141; 6 literal terms
\begin{align*}
q_{141} &={}-x_{14}^{2} +2\,x_{14}x_{15} +x_{14}x_{66} -x_{15}^{2}\\
&\quad -x_{15}x_{66} +x_{189}
\end{align*}
% residual 142; 4 literal terms
\begin{align*}
q_{142} &={}x_{14}x_{189} -x_{15}x_{189} -x_{67}x_{189} +x_{190}
\end{align*}
% residual 143; 4 literal terms
\begin{align*}
q_{143} &={}x_{14}x_{190} -x_{15}x_{190} -x_{68}x_{190} +x_{191}
\end{align*}
% residual 144; 4 literal terms
\begin{align*}
q_{144} &={}x_{14}x_{191} -x_{15}x_{191} -x_{69}x_{191} +x_{192}
\end{align*}
% residual 145; 4 literal terms
\begin{align*}
q_{145} &={}x_{14}x_{192} -x_{15}x_{192} -x_{70}x_{192} +x_{193}
\end{align*}
% residual 146; 4 literal terms
\begin{align*}
q_{146} &={}x_{14}x_{193} -x_{15}x_{193} -x_{71}x_{193} +x_{194}
\end{align*}
% residual 147; 4 literal terms
\begin{align*}
q_{147} &={}x_{14}x_{194} -x_{15}x_{194} -x_{72}x_{194} +x_{195}
\end{align*}
% residual 148; 4 literal terms
\begin{align*}
q_{148} &={}x_{14}x_{195} -x_{15}x_{195} -x_{73}x_{195} +x_{196}
\end{align*}
% residual 149; 2 literal terms
\begin{align*}
q_{149} &={}-x_{79}x_{196} +x_{197}
\end{align*}
% residual 150; 2 literal terms
\begin{align*}
q_{150} &={}-x_{81}x_{197} +x_{198}
\end{align*}
% residual 151; 5 literal terms
\begin{align*}
q_{151} &={}-x_{10}x_{11} +x_{10}x_{12} +x_{11}x_{12} -x_{12}^{2}\\
&\quad +x_{199}
\end{align*}
% residual 152; 4 literal terms
\begin{align*}
q_{152} &={}x_{13}x_{198} +2\,x_{198} -x_{198}x_{199} +x_{200}
\end{align*}
% residual 153; 4 literal terms
\begin{align*}
q_{153} &={}-x_{13}x_{200} -x_{199}x_{200} -2\,x_{200} +x_{201}
\end{align*}
% residual 154; 4 literal terms
\begin{align*}
q_{154} &={}x_{14}x_{198} -x_{15}x_{198} -x_{74}x_{198} +x_{202}
\end{align*}
\endgroup

% END embedded residuals203.tex


\clearpage
\phantomsection
\addcontentsline{toc}{section}{References}
\begin{thebibliography}{99}
\small\raggedright
\bibitem{p163} ProveIt, package 163, \emph{Lower degree and fewer witnesses
for Diophantine consistency}, 25 October 2026. In particular
\code{compact/PROOF.md}, sections 3--9.
\url{https://github.com/VladimirReshetnikov/ProveIt/tree/fb1bed5eb7196cf35abfaa746c18c6ee7eae6cfb/Computability/DiophantineConsistency/con-zfc-diophantine-forms/163-trace-radical}.
\bibitem{p164} ProveIt, package 164, \emph{Unsquared factors and linear
selectors for Diophantine consistency}, 25 October 2026. The positive
inverse-Lucas proof and the natural canonical CRT proof are the immediate
arithmetic parents.
\url{https://github.com/VladimirReshetnikov/ProveIt/tree/fb1bed5eb7196cf35abfaa746c18c6ee7eae6cfb/Computability/DiophantineConsistency/con-zfc-diophantine-forms/164-unsquared-factors-linear-selectors}.
\bibitem{p159} ProveIt, package 159, \emph{A direct quartic for Con(ZFC)
in 234 integer variables}, 25 October 2026. Common-Hessian projection
and independent algebraic certificates.
\url{https://github.com/VladimirReshetnikov/ProveIt/tree/fb1bed5eb7196cf35abfaa746c18c6ee7eae6cfb/Computability/DiophantineConsistency/con-zfc-diophantine-forms/159-quadratic-projection-234}.
\bibitem{p156} ProveIt, package 156, \emph{Integral projection lattice},
\code{research_notes/EXPERIMENT_REVIEW.md}. Prior measured seed/vacuity split.
\url{https://github.com/VladimirReshetnikov/ProveIt/blob/fb1bed5eb7196cf35abfaa746c18c6ee7eae6cfb/Computability/DiophantineConsistency/con-zfc-diophantine-forms/156-integral-projection-lattice/research_notes/EXPERIMENT_REVIEW.md}.
\bibitem{p152} ProveIt, package 152, \emph{Lucas corrections and elimination
of arithmetic witnesses}, 25 October 2026. The earlier forward Lucas
correction, retained in the source closure.
\url{https://github.com/VladimirReshetnikov/ProveIt/tree/fb1bed5eb7196cf35abfaa746c18c6ee7eae6cfb/Computability/DiophantineConsistency/con-zfc-diophantine-forms/152-lucas-corrections}.
\bibitem{signedcenters} ProveIt, packages 153, 154 and 165. Earlier
whole-expression polynomial centers for eleven unrestricted integer variables:
\path{153-nearest-factorial/signed/PROOF.md}, section 6;
\path{154-normalized-residues-round19/report/scalar.tex}; and
\path{165-wider-budgets-round20/scalar/PROOF.md}.
\url{https://github.com/VladimirReshetnikov/ProveIt/blob/fb1bed5eb7196cf35abfaa746c18c6ee7eae6cfb/Computability/DiophantineConsistency/con-zfc-diophantine-forms/165-wider-budgets-round20/scalar/PROOF.md}.
\bibitem{kendpoint} ProveIt, \code{KEndpoint.lean},
\code{KSeedDirect222.lean}, \code{ConstructibleChoice.lean}.
Named C222/ZF/ZFC interfaces, with the exact selected source copies and
commit identifiers recorded in the artifact ledgers.
\bibitem{jones} J. P. Jones, \emph{Universal Diophantine Equation},
Journal of Symbolic Logic 47 (1982), 549--571.
\url{https://doi.org/10.2307/2273588}. Background for comparison only;
the direct quartic component does not use this universal equation.
\end{thebibliography}
\end{document}
